% 本文件由 LaTeX 排版 Agent 根据有效 translation.json 自动生成。
% author=Justin Martyr; work_id=0126; title=第一护教篇

\chapter{第一护教篇}

% structure: p0001 source=h1 target=omitted reason=page-title-already-rendered-by-parent

% structure: p0002 source=h2 target=section reason=relative-heading-rank; base=h2

\section{第一章 致辞}

致皇帝\Person{提图斯·埃利乌斯·哈德良·安东尼·庇护·奥古斯都·凯撒}，及其子、哲学家\Person{维里西穆斯}，以及凯撒的亲子、庇护的养子、热爱学问的哲学家\Person{卢修斯}，并神圣的元老院及全体罗马人民：我，\Person{尤斯定}，\Person{普里斯库斯}之子、\Person{巴基乌斯}之孙，原籍\Place{巴勒斯坦}的\Place{弗拉维亚·尼阿波利斯}，现代表所有不公正地遭人仇恨、肆意凌辱的各民族之人，呈上这篇致辞与请愿书，我自己亦是其中之一。\par

% structure: p0004 source=h2 target=section reason=relative-heading-rank; base=h2

\section{第二章 要求公正}

理性引导真正虔诚而热爱哲学的人，只尊崇并爱慕真理，拒绝追随传统的见解，若这些见解毫无价值。因为不仅健全的理性使我们拒绝接受那些言行有误者的指引，而且热爱真理的人，即使面对死亡的威胁，也当选择行事说话合乎正义，甚至置生死于度外。那么，你们既被称为虔诚者与哲学家，是正义的护卫者、学问的爱好者，就请留心听我的致辞；若你们果真如此，便会显明出来。我们前来，并非藉此文阿谀奉承，也非以言词取悦，而是恳请你们经过准确而深入的查究之后，作出判决，不要被偏见或取悦迷信者的欲望所蒙蔽，也不要被非理性的冲动或长期流传的恶言所驱使，而作出终究对你们自己不利的决定。因为我们相信，除非我们被证明为恶徒或坏人，否则无人能加害于我们；而你们，虽能杀死我们，却不能伤害我们。\par

% structure: p0006 source=h2 target=section reason=relative-heading-rank; base=h2

\section{第三章 要求司法调查}

但恐怕有人认为这是无理而鲁莽的言辞，我们要求调查对基督徒的指控；若这些指控属实，就应依其罪受罚；〔或者，更确切地说，我们自己也愿意惩罚他们。〕但若无人能证明我们有罪，真正的理性禁止你们因恶意的谣言而冤枉无辜之人，甚至更冤枉你们自己，因为你们不以判断，而以激情来治理事务。每个头脑清醒的人都会宣称，唯一公平公正的调整是：臣民对自己的生活与教义提出无可挑剔的说明，而统治者则当服从于虔诚与哲学，而非暴力与暴政来作出判决。这样，统治者和被统治者都能获益。因为古代某位哲人曾说过：\Quote{除非统治者和被统治者都成为哲学家，否则国家不可能幸福。}因此，我们的责任是让所有人都有机会审视我们的生活与教导，以免我们因那些习惯于无知我们之事的人，而遭受他们因心灵盲目所应得的惩罚；而你们的工作则是，在听取我们的陈述后，按理性要求成为公正的法官。因为若你们得知真相后仍不秉公行事，你们在天主面前将无可推诿。\par

% structure: p0008 source=h2 target=section reason=relative-heading-rank; base=h2

\section{第四章 基督徒仅仅因名而被不公正地定罪}

仅仅应用一个名称，若不结合该名称所蕴含的行为，并不能决定善恶；事实上，至少从我们被指控的这个名称来判断，我们是最优秀的人。但我们认为，若我们被证实为恶徒，则不应因这个名称而请求无罪开释；反之，若我们无论在称呼自己的方式上，还是在作为公民的行为上，都没有犯下任何过错，那么你们就有责任极力避免因不公正地惩罚未被定罪的人而招致公正的惩罚。因为一个名称既不能产生赞美，也不应导致惩罚，除非在行为上证明有善或恶。而你们在指控你们自己人中，不会在定罪之前就惩罚他们；但对我们，你们却把名称当作证据来定罪，尽管就名称而言，你们更应惩罚我们的指控者。因为我们被指控为基督徒，而仇恨\Emph{优秀}（\OriginalTerm{优秀}{Chrestian}）是不义的。再者，若被指控者中有人否认这个名称，说自己不是基督徒，你们就释放他，因为没有不利于他的证据；但若有人承认自己是基督徒，你们就因这个承认而惩罚他。公正要求你们既审问承认者，也审问否认者，通过他们的行为，可以显明各人是何种人。因为正如一些受导师\Person{基督}教导而不否认他的人，在被质问时给别人以鼓励，同样，那些行为邪恶的人也很可能给那些不假思索就指控所有基督徒不虔不义的人提供口实。这也是不对的。因为在哲学中，也有人空有其名与外表，却行不出与宣称相称的事；而且你们很清楚，古代那些见解与教导极为不同的人，却都被统称为哲学家。其中一些教人无神论；而你们中的诗人以\Person{朱庇特}与自己儿女的污秽之事取乐。如今那些采纳此类教诲的人并未受你们约束；相反，你们将奖品与荣誉授予那些以悦耳之音侮辱神明的人。\par

% structure: p0010 source=h2 target=section reason=relative-heading-rank; base=h2

\section{第五章 基督徒被指控无神论}

为何如此？在我们这些承诺不作恶、也不持这些无神论观点的人身上，你们不审查对我们的指控，反而屈从于无理性的激情和恶鬼的煽动，不加考虑、不按判断就惩罚我们。因为必须说出真理：自古以来，这些恶鬼显现自身，污辱妇女、败坏儿童，并向人展示如此恐怖的景象，以致那些不用理性判断所行之事的人被恐惧所震慑；且被恐惧所挟，不知这些是鬼，便称它们为神，给每个鬼各自选择的名称命名。当\Person{苏格拉底}试图以真理性与考察揭露这些事、将人从鬼中解救出来时，那些鬼自身便借助喜悦不义的人，以无神论者和亵渎者的罪名处死了他，指控他\Quote{引进新神祇}；而在我们身上，他们施展了类似的伎俩。因为不仅是在希腊人中，\OriginalTerm{理性}{Logos}通过\Person{苏格拉底}得胜而谴责了这些事，在野蛮人中，\OriginalTerm{圣言}{Logos}亲自取了形体、成为人、被称为\Person{耶稣基督}，也谴责了它们；我们服从祂，不仅否认那些行如此之事者为神，而且断言他们是邪恶不敬的鬼，其行为甚至无法与向往德行的人相比。\par

% structure: p0012 source=h2 target=section reason=relative-heading-rank; base=h2

\section{第六章 \OriginalTerm{无神论}{atheism}的指控被驳斥}

因此我们被称为无神论者。我们承认，就这类神祇而言，我们是无神论者；但关于至真天主、公义和节制及其他德行的圣父（祂远离一切不洁），我们不仅敬拜祂，也敬拜从祂而来、教导我们这些事的圣子，以及跟随祂、效法祂的其他善天使众军，还有圣神；我们用理性和真理认识他们，并毫无保留地按所受的教导向每个愿意学习的人宣讲。\par

% structure: p0014 source=h2 target=section reason=relative-heading-rank; base=h2

\section{第七章 每位基督徒应按其生活受审}

但有人会说：此前有些人已被逮捕并定罪为恶人。因为你们多次在分别审查每个被告的生活后进行判决，但并非由于我们所谈论的那些人。这一点我们承认：如同在希腊人中，所有教授自己所喜理论的人都被通称为\Quote{哲学家}，尽管他们的学说各异；同样，在野蛮人中，这个招致众多指控的名称是那些真正有智慧和貌似有智慧者所共有的。因为所有人都被称为基督徒。因此我们要求，凡被指控到你们面前的人，其行为都应受审判，以便每个被定罪的人都作为恶人受惩罚，而非作为基督徒；若某人显然清白，则他应被释放，因为仅凭他是基督徒这一事实，他并未行不义。我们不会要求你们惩罚我们的控告者，他们已被自己的邪恶和对正义的无知所充分惩罚。\par

% structure: p0016 source=h2 target=section reason=relative-heading-rank; base=h2

\section{第八章 基督徒宣认他们对天主的信仰}

请相信我们讲这些是为你们着想；因为当我们受审时，我们本可以否认自己是基督徒，但我们不愿靠说谎而活。因为被对永恒纯洁生命的渴望所推动，我们追求与万物的圣父和造物主同在的居所，并急于宣认我们的信仰，因为我们确信并深信：那些通过行为向天主证明自己跟随祂、喜爱与祂同在（那里没有引起扰乱的罪恶）的人，必能获得这些。简言之，这就是我们从基督所期待、所学习并教导的。同样，\Person{柏拉图}惯常说\Person{剌达曼堤斯}和\Person{米诺斯}将惩罚来到他们面前的恶人；而我们也说同样的事将发生，但藉着基督之手，恶人在同他们灵魂重新结合的身体中（这身体现在将受永罚），且不止如\Person{柏拉图}所说的仅一千年。若有人说这不可信或不可能，那么我们这错误只关系到我们自己，不涉及任何人，只要你们不能证明我们行了任何损害。\par

% structure: p0018 source=h2 target=section reason=relative-heading-rank; base=h2

\section{第九章 崇拜偶像的愚蠢}

我们也不以许多祭献和花环来尊崇那些被人类造出、置于神龛并称为神祇的偶像；因为我们看见这些偶像无灵魂、已死亡，没有天主的形状（因我们认为天主不具有某些人所说模仿以尊崇祂的形状），却拥有那些显现的恶鬼的名字和形状。何必向早已知道的你们叙述，工匠们\BibleRef{依 44:9}如何雕刻、切割、铸造、锤打材料，使之成形？他们常从卑贱的器皿出发，仅改变形状，制作所需样式的像，就称之为神；我们认为这不但无知，更是对天主的侮辱，因天主拥有不可言喻的荣耀和形状，却让祂的名字附着于可朽坏、需不断服侍之物上。这些工匠既放纵无度，又（不细说）精通各种恶行，你们很清楚；他们甚至败坏与自己一同做工的女子。何其愚昧！竟说放荡之人能为你们的敬拜塑造并制造神祇，而你们又任命这样的人为供奉神祇的庙宇守护者；却未认识到，连想或说人是神的守护者都是不合法的。\par

% structure: p0020 source=h2 target=section reason=relative-heading-rank; base=h2

\section{第十章 应如何事奉天主}

我们由传统得知，天主并不需要人们所能提供的物质祭品，因为祂本身就是万物的供给者。我们受教导、深信并相信，祂只接纳那些效法祂内所存美善的人，即节制、正义、仁爱，以及一切属于那位没有适当名称之天主的德性。我们又受教导，祂在起初出于自己的良善，为了人的缘故，从无定形之物创造了万物；若人们以行为表明自己配得上祂这计划，他们就被认为配得——如我们所领受的——与祂一同为王，脱离朽坏与痛苦。因为，正如起初祂在我们不存在时创造了我们，同样，我们认为，那些选择祂所喜悦之事的人，因着他们的选择，被认为是配得不朽坏并与祂共融的。因为最初的存在并不在我们自己的能力之内；而为了使我们能追随祂所喜悦的事，藉由祂自己所赋予我们的理性官能选择它们，祂既劝勉我们，也引导我们走向信德。我们认为，所有人不受到阻止去学习这些事，反而被催促去学习，对所有人都是有益的。因为人类法律无法达成的约束，圣言——既然祂是天主——本来能够成就，若不是邪恶的魔鬼，以每人心中那邪恶的私欲（它纷繁地引人走向各种恶习）为盟友，散布了许多虚假而亵渎的指控，其中无一与我们有关。\par

% structure: p0022 source=h2 target=section reason=relative-heading-rank; base=h2

\section{第十一章 基督徒所盼望的王国}

当你们听说我们盼望一个王国时，你们不加探究就以为我们说的是人间王国；但我们所说的，是与天主同在的王国，这从那些被指控为基督徒之人的信仰宣认中也可看出，尽管他们知道死亡是加给如此宣认者的惩罚。因为如果我们盼望的是人间王国，我们也会否认我们的基督，以免被杀；并且我们会竭力逃避侦查，好得到我们所期望的。但既然我们的思念不专注于现世，当人们断绝我们时，我们也不以为意；因为死亡也是一笔无论如何都要偿还的债。\par

% structure: p0024 source=h2 target=section reason=relative-heading-rank; base=h2

\section{第十二章 基督徒生活在天主的眼下}

我们在促进和平时，比所有人都更是你们的帮手和盟友，因为我们持守这观点：无论是恶人、贪婪者、阴谋家，还是善人，都不可能逃过天主的注意，并且各人按照自己行为的价值，走向永远的惩罚或救恩。因为如果所有人都知道这一点，没有人会哪怕短暂地选择邪恶，因为他知道自己将走向永远的火刑；反而会竭力约束自己，以美德装饰自己，好得到天主的美善恩赐，并逃脱惩罚。因为那些因你们所施行的法律和惩罚而试图在犯罪时逃避侦查的人（他们犯罪时也以为完全可能逃过你们的侦查，因为你们不过是人），这些人若得知并深信没有什么事——无论是实际做的还是仅属意图——能逃过天主的认知，他们就会因所威胁的惩罚而竭力正当地生活，连你们自己也必得承认。但你们似乎害怕所有人成为义人，你们就没有人可惩罚了。这会是刽子手们的关切，而不是好君王的关切。但正如我们先前所说，我们深信这些事是由邪恶的灵所促成的，他们甚至向那些生活不合理的人要求祭献和事奉；但对你们，我们假定你们追求虔诚和哲学之名，不会做任何不合理的事。但如果你们也像愚昧人一样，宁要习俗而不要真理，那么你们就做你们有能力做的事吧。但那些重视意见过于真理的统治者，他们的能力就像荒漠中的强盗一样。而你们不会成功，这是圣言所宣告的；在生了祂的天主之后，我们知道没有比祂更王权、更公正的统治者了。因为正如人人都避免继承父亲的贫穷、痛苦或卑微一样，凡是圣言禁止我们选择的，明智的人都不会选择。我说，这一切事情的发生，是我们的导师所预言的，祂就是万有之父和统治者的天主之子、宗徒，耶稣基督；我们也从祂得了基督徒的名。因此，我们对祂所教导的一切事更加确信，因为祂事先预言将要发生的一切，如今正看到在发生；这是天主的工作：在事情发生之前就说出来，并且正如所预言的，显明它正在发生。我们本可以就此打住，不再多说，认为我们所要求的是公正和真实的；但因我们深知，突然改变一个被无知占据的心并不容易，我们打算再补充几点，为说服那些爱真理的人，知道用真理驱散无知并非不可能。\par

% structure: p0026 source=h2 target=section reason=relative-heading-rank; base=h2

\section{第十三章 基督徒理性地事奉天主}

那麼，還有哪個清醒的人會不承認我們並非無神論者呢？因為我們敬拜這個宇宙的造物主，並按照所受的教導宣稱，祂並不需要血流、奠酒和馨香；我們以祈禱和感恩來盡力讚美祂，感謝祂所供應我們的一切，因為我們所受的教導是：唯一配得上祂的尊榮，不是用火燒毀祂為我們養生所創造之物，而是將它們用於我們自己和有需要的人，並懷著感激之情，透過祈禱和詩歌為我們的受造、為一切健康的途徑、為各類事物不同的性質、為季節的變換而獻上感謝；並且在祂面前呈上請求，使我們因著對祂的信德而得以在不朽中重生。教導我們這一切的是\Person{耶稣基督}，祂正是為此而降生，並在\Person{提庇留·凯撒}時代，於\Person{般雀·比拉多}——猶太行省總督——任內被釘在十字架上；而我們合理地敬拜祂，因為我們知道祂就是那真實天主自己的\OriginalTerm{聖子}{Son}，並將祂置於第二位，將先知之神置於第三位，這我們將會證明。因為他們指責我們的瘋狂在於：我們將一個被釘十字架的人置於永恆不變、創造萬有的天主之下第二位；但他們沒有辨明其中所隱藏的奧秘，而我們向你們解釋這奧秘時，請求你們留意。\par

% structure: p0028 source=h2 target=section reason=relative-heading-rank; base=h2

\section{第十四章：邪魔歪曲基督徒的教義}

因為我們預先警告你們要提防，免得我們所指控的那些邪魔欺騙你們，使你們完全偏離閱讀和理解我們所說的話。因為它們竭力要使你們成為它們的奴隸和僕人；有時藉著夢中的顯現，有時藉著魔術的欺騙，它們制伏了一切沒有為自己的救恩而作出堅決抵抗的人。因此，我們自從被\OriginalTerm{聖言}{Word}說服之後，也遠離它們（即邪魔），並藉著祂的\OriginalTerm{聖子}{Son}跟隨那唯一不受生的天主——我們從前縱情於淫亂，如今卻唯獨擁抱貞潔；我們從前使用邪術，如今卻將自己奉獻給善的、不受生的天主；我們從前視財富和產業的獲得高於一切，如今卻將所有的東西歸於公用，並分給每一個有需要的人；我們從前彼此仇恨和殘殺，且因不同的習俗不願與不同部族的人同住，如今，自從\Person{基督}來臨之後，卻與他們親切相處，並為我們的仇敵祈禱，努力勸說那些無理仇恨我們的人，使他們遵照\Person{基督}的善訓生活，好使他們能與我們一同分享來自萬有統治者天主的同一喜樂的希望報酬。但為了避免我們顯得是在詭辯，我們認為在向你們提供所應許的解釋之前，引用\Person{基督}親自頒布的幾條教訓是合適的。而你們，作為有權柄的統治者，有責任查問我們是否真實地接受了這些教導並如此教導別人。祂說出的話語簡短而精煉，因為祂不是詭辯家，但祂的話語是天主的德能。\par

% structure: p0030 source=h2 target=section reason=relative-heading-rank; base=h2

\section{第十五章：基督親自教導了什麼}

论及贞洁，祂曾如此说：\Quote{凡注视妇女而贪恋她的，已经在心里与她犯了奸淫。} 又说：\Quote{如果你的右眼使你跌倒，就把它剜出来；因为带着一只眼进入天国，比有两只眼而被投入永火更好。} 又说：\Quote{凡娶那被丈夫休弃的妇人的，就是犯奸淫。} 又说：\Quote{有些人是生来阉人，有些人是被人阉的，有些人是为天国而自阉的；但这话不是人人都能领受的。} \BibleRef{玛 19:12} 因此，凡按人间法律再婚的，在我们主眼中都是罪人，凡注视妇女而贪恋她的也是如此。因为不仅实际犯奸淫的为祂所弃绝，连贪恋犯奸淫的亦然：因为不仅我们的行为，连我们的思念，在天主面前也是显明的。许多自幼便作基督门徒的男女，到六十或七十岁仍保持纯洁；我敢夸口，能从每一民族中举出这样的人。至于那无数改过自新、学习这些道理的人，我又何须多说？因为基督来，不是召义人，也不是召贞洁者悔改，而是召不虔敬、纵欲、不义的人；祂的话是：\Quote{我不是来召义人，而是召罪人悔改。} \BibleRef{玛 9:13} 因为天上的父更愿罪人悔改，而非受罚。论到我们对众人的爱，祂这样教导：\Quote{如果你们爱那爱你们的人，你们做了什么特别的呢？就是淫乱的人也这样做。但我告诉你们：要为你们的仇敌祈祷，爱那恨你们的人，祝福那诅咒你们的人，为那凌辱你们的人祈祷。} \BibleRef{玛 5:46} 至于我们应当周济穷人，不为荣耀行事，祂说：\Quote{凡求你的，就给他；向你借贷的，不要转身走开。因为你们若借给那些你们希望偿还的人，你们做了什么特别的事呢？就连税吏也这样做。不要为自己积蓄财宝在地上，地上有虫蛀、锈蚀，也有贼挖洞偷窃；但要在天上积蓄财宝，那里没有虫蛀、锈蚀，也没有贼挖洞偷窃。因为人若赚得全世界，却丧失自己的灵魂，有什么益处呢？或者人还能拿什么交换自己的灵魂？所以，要把财宝积蓄在天上，那里没有虫蛀、锈蚀。} 又说：\Quote{你们应当慈悲，如同你们的父是慈悲的一样，祂使太阳升起，光照恶人，也光照善人；降雨给义人，也给不义的人。不要忧虑吃什么、穿什么：你们不比飞鸟和走兽更贵重吗？天主尚且养活它们。所以，不要忧虑吃什么、穿什么，因为你们的天父知道你们需要这一切。你们先该寻求天国，这一切自会加给你们。因为你们的财宝在哪里，你们的心也在那里。} 又说：\Quote{不要为了让人看见而做这些事，否则你们在天父那里就没有赏报了。} \BibleRef{玛 6:1}\par

% structure: p0032 source=h2 target=section reason=relative-heading-rank; base=h2

\section{第十六章　论忍耐与起誓}

论到我们当含忍冤屈、乐意服事众人、不轻易发怒，祂这样说：\Quote{有人打你的右脸，连左脸也转过来给他；有人拿你的外衣或里衣，不要阻止他。凡向你发怒的，难免受火的刑罚。有人强迫你走一里路，你就同他走两里。你们的光也当这样照在人前，叫他们看见你们的好行为，便将荣耀归给你们在天上的父。} 因为我们不该争斗；祂也不愿我们效法恶人，却劝我们用忍耐和温柔，引导众人脱离羞耻和爱恶。这事实可由许多先前与你们想法相同的人身上得到证明，他们因目睹邻人生活中的恒心，或因看到同伴被欺压时的非常容忍，或因与同交易者的诚实，而改变了暴戾专横的性情。\par

关于我们完全不发誓、常说实话，祂吩咐如下：\Quote{你们什么誓都不可起；你们的话，是就说是，非就说非；此外多说的，便是出于邪恶。} \BibleRef{玛 5:34} 关于我们应当唯独崇拜天主，祂这样劝导：\Quote{最大的诫命是：你要崇拜上主你的天主，唯独事奉祂，全心、全力爱那造你的上主天主。} \BibleRef{谷 12:30} 当一个人来对祂说：\Quote{善师，} 祂回答说：\Quote{除了那创造万有的独一天主外，没有一个是善的。} \BibleRef{玛 19:6} 凡不按祂的教导生活的人，即使口中承认基督的诫命，也应明白他们不是基督徒；因为不是那口称的，而是那实行的人才得救，按祂的话：\Quote{不是凡对我说\InnerQuote{主啊，主啊}的人，都能进入天国，而是那承行我在天之父旨意的人。因为凡听了我这些话而实行的，就是听从那派遣我来的。许多人要对我说：\InnerQuote{主啊，主啊，我们不是因你的名吃过喝过，行过奇事吗？}那时我必要对他们说：\InnerQuote{离开我，你们这些作恶的人！}那时，义人将发光如同太阳，恶人被投入永火，那里有哀号和切齿。因为许多人要假借我的名来，外面披着羊皮，里面却是凶残的豺狼。你们可凭他们的果实辨别他们。凡不结好果子的树，都要砍倒投入火中。} 凡不按这些教训生活、徒有基督徒之名的人，我们要求你们惩处他们。\par

% structure: p0035 source=h2 target=section reason=relative-heading-rank; base=h2

\section{第十七章　基督教导服从政权}

我们比任何人都更愿意向你们所指定的人缴纳常税和额外税，这是我们从祂所学的；因为那时有些人来到祂面前，问祂是否该纳税给\OriginalTerm{凯撒}{Cæsar}，祂回答说：\Quote{告诉我，这钱币上是谁的像？}他们说：\Quote{\OriginalTerm{凯撒}{Cæsar}的。}祂又对他们说：\Quote{\OriginalTerm{凯撒}{Cæsar}的归\OriginalTerm{凯撒}{Cæsar}，天主的归天主。}因此，我们只朝拜天主，但在其他事上，我们乐意服事你们，承认你们是人类的君王和统治者，并祈求你们在王权之下也拥有健全的判断。但如果你们不理会我们的祈祷和坦诚的解释，我们不会受损，因为我们相信（其实是确信）每个人都要按自己行为的功过在永火中受罚，并要按他从天主领受的能力交账，正如基督所说：\BibleQuote{多给谁，就向谁多要。}\BibleRef{路 12:48}\par

% structure: p0037 source=h2 target=section reason=relative-heading-rank; base=h2

\section{第十八章 不死与复活的证明}

因为请回想每位先前君王的结局，他们都是如何同归一死；这死亡若是归于无知觉，那对所有恶人倒是好事。但既然所有曾活过的人都保有感觉，且永罚已经为他们（即恶人）存留，你们要小心，不要忽略被说服，并相信这些事是真实的。因为即使是招魂术、你们用纯洁孩童所行的占卜、召唤已离世的人魂、以及那些在术士中称为托梦者和助灵（\OriginalTerm{家神}{Familiars}）的事，连同那些精通此道者所做的一切——就让这些说服你们：即使在死后，灵魂仍有知觉；还有那些被死人之灵抓住并甩来甩去的人（你们称为附魔者或疯子）；以及你们视为神谕的，无论是\OriginalTerm{安菲洛科斯}{Amphilochus}、\OriginalTerm{多多纳}{Dodana}、\OriginalTerm{皮托}{Pytho}，以及其他现有的这类神谕；还有你们作家的见解，如\OriginalTerm{恩培多克勒}{Empedocles}、\OriginalTerm{毕达哥拉斯}{Pythagoras}、\OriginalTerm{柏拉图}{Plato}、\OriginalTerm{苏格拉底}{Socrates}、\OriginalTerm{荷马}{Homer}的深渊，\OriginalTerm{尤利西斯}{Ulysses}下到阴间察看这些事，以及一切同类的话语。你们既然赐予这些事如此恩惠，也请赐予我们，我们并不比他们更少，而是更坚定地相信天主；因为我们期待再次得回自己的躯体，即使它们已死并被埋入土中，因为我们坚信：在天主没有不可能的事。\par

% structure: p0039 source=h2 target=section reason=relative-heading-rank; base=h2

\section{第十九章 复活是可能的}

对于任何有思想的人来说，还有什么比以下更不可思议呢：如果我们不在身体内，有人却说从一小滴人种可以长出骨头、筋和肉，形成我们所见的形状？现在就假设地说：如果你们自己不是现在这样，也不是由这样的父母[和原因]所生，有人向你们展示人种和一个人的图画，并确信地说从这样的物质能产生这样的存在，你们在亲见产生之前会相信吗？没有人敢否认[这样的说法会超越信仰]。同样，你们现在不信，是因为从未见过死人复活。但正如起初你们不会相信从那一小滴能产生这样的人，而现在你们看见他们如此产生，同样，也要判断：人的身体在分解后，如同种子被化解入土，在天主所定的时候复活并穿上\OriginalTerm{不朽}{incorruption}，并不是不可能的。因为那些说万物回归其原质，且除此之外连天主自己也不能做什么的人，他们所想象的天主的能力，我们无法理解；但这一点我们看得很清楚：他们不会相信他们自己能够从那样的材料成为现在所见的自己和整个世界，这是可能的。而且，相信甚至在我们本性和人看来不可能的事，比像世界其余的人那样不信更好，这是我们所学的；因为我们知道我们的主耶稣基督说过：\BibleQuote{在人不可能，在天主却可能，}\BibleRef{玛 19:26}又说：\BibleQuote{不要害怕那些杀害你们，此后不能再做什么的人；但要害怕那死后能连灵魂带身体投进地狱的那一位。}\BibleRef{玛 10:28}。地狱是一个地方，在那里那些生活邪恶、并且不相信天主藉基督所教导我们的事将要成就的人，将受惩罚。\par

% structure: p0041 source=h2 target=section reason=relative-heading-rank; base=h2

\section{第二十章 异教徒与基督教义的类比}

\OriginalTerm{西彼拉}{Sibyl}和\OriginalTerm{希斯塔斯普}{Hystaspes}说过，天主将毁灭可朽之物。而称为\OriginalTerm{斯多葛派}{Stoics}的哲学家们教导说，连天主自己也将被分解为火，他们还说世界要藉这轮转重新形成；但我们理解，万物的创造者天主，超越那些将要改变的事物。因此，如果我们在某些点上与你们所尊崇的诗人和哲学家教导相同，在其他点上我们的教导更丰富、更神圣，并且只有我们为我们所断言的事提供证明，为什么我们所受的仇恨比所有其他人更不公义呢？因为我们说万物由天主产生并安排成一个世界，我们似乎在说出\OriginalTerm{柏拉图}{Plato}的学说；我们说万物将被焚烧，似乎是在说斯多葛派的学说；我们肯定恶人的灵魂在死后仍具有感觉而受惩罚，好人的灵魂脱离惩罚度过有福的存在，我们似乎在说出诗人和哲学家相同的话；我们坚持人不应该崇拜自己手所做的工，我们说的恰恰是喜剧诗人\OriginalTerm{米南德}{Menander}和其他类似作家所说过的，因为他们已经宣称：工人大于作品。\par

% structure: p0043 source=h2 target=section reason=relative-heading-rank; base=h2

\section{第二十一章 基督历史的类比}

当我们又说圣言——祂是天主首生者——无两性结合而受生，且祂，耶稣基督，我们的导师，被钉十字架而死，又复活，升了天堂，我们所说的与你们所信的关于那些你们尊为朱庇特之子的，并无不同。因为你们知道，你们所尊重的作家将多少儿子归于朱庇特：墨丘利，众神的消息者与导师；埃斯科拉庇俄斯，虽是一位大医师，却被雷电击打，如此升天；巴克斯也同样，被肢解后；赫拉克勒斯，当他投身火焰以逃避劳苦时；勒达与狄奥斯库里之子；达那厄之子珀耳修斯；以及贝勒罗丰，虽生于凡人，却乘珀伽索斯马升天。至于阿里阿德涅，以及像她一样被宣布位列星辰的，我还能说什么呢？还有你们中间去世的皇帝，你们认为他们堪受神化，并为他们的缘故，有人发誓说看见焚烧的凯撒从火葬堆升天。关于这些所谓的朱庇特之子，每人所记载的作为，对于已知者无需赘述。只需说：这些是为年轻学子的利益和鼓励而写的；因为所有人都以效法神祇为荣。但每个有良好教养的灵魂绝不应如此思想神祇，竟相信朱庇特本人——万有的治理者与创造者——既是一位弑父者又是弑父者之子，且被卑劣可耻的快乐之爱征服，进入伽倪墨得斯和许多他所侵犯的女人那里，而他的儿子们也做类似的事。但是，正如我们上文所说，是邪恶的魔鬼行了这些事。而我们学到，只有那些在圣洁和美德中亲近天主的人才被神化；我们相信，那些邪恶生活而不悔改的人将在永恒之火中受罚。\par

% structure: p0045 source=h2 target=section reason=relative-heading-rank; base=h2

\section{第二十二章　论基督圣子身份的类比}

再者，天主子称为耶稣，即使仅按普通生育为人，但因祂的智慧，仍堪称为天主子；因为所有作者都称天主为人和神祇之父。如果我们断言天主的圣言以一种特殊方式、不同于普通生育由天主所生，那么，如上所述，这对你们而言并非奇异之事，你们也说墨丘利是天主的天使圣言。但若有人反对祂被钉十字架，在这件事上祂也与你们那些传说中像我们现在列举的受苦的所谓朱庇特之子相等。因为他们死时的苦难被记载为并非全都相同，而是多样的；因此，即使从祂受苦的特殊性来看，祂似乎也不比他们差；相反，正如我们在这篇论述的前面部分所承诺的，我们现在将证明祂更优越——或者已经证明了——因为优越者由祂的行动显明。如果我们甚至肯定祂由童贞女所生，请将此与你们所接受的关于珀耳修斯的事同等看待。而当我们说祂使瘸子、瘫痪者和生来瞎眼的人痊愈时，我们所说的似乎与所传埃斯科拉庇俄斯所行的事非常相似。\par

% structure: p0047 source=h2 target=section reason=relative-heading-rank; base=h2

\section{第二十三章　论证}

为使这一点现在对你们显明——（首先）我们所主张的合乎基督及祂之前的先知所教导我们的一切，唯有这些是真实的，并且比所有已存在的作家都更古老；我们要求被认可，并非因我们说了与这些作家相同的话，而是因我们说了真话；（其次）耶稣基督是唯一天主所生的适格之子，祂是圣言和首生者及德能；并按照祂的旨意成为人，祂教导我们这些事，为使人性归正和恢复；（第三）在祂成为人、居于人中间之前，有些受前述魔鬼影响的人，通过诗人预先讲述了那些作为实际发生的事，而他们虚构编造并叙述了这些事，正如他们也使那关于我们的无耻和不虔敬行为的诽谤报告被捏造出来，这些行为既无见证也无证据——我们将提出以下证明。\par

% structure: p0049 source=h2 target=section reason=relative-heading-rank; base=h2

\section{第二十四章　异教崇拜的多样性}

首先[我们提出证明]，因为，虽然我们说了与希腊人相似的话，却唯独因基督这名而遭憎恨，尽管我们没有行恶，却被作为罪人处死；而别处的其他人则崇拜树木、河流、老鼠、猫、鳄鱼及许多无理性的动物。同一动物也非被所有人尊崇；而是在一处崇拜这个，在另一处崇拜那个，因此，在他们彼此眼中，所有人都因不崇拜同一对象而成为亵渎者。而这是你们对我们的唯一控告：我们不敬拜你们所敬拜的同样的神祇，也不向死人献奠酒和脂油之香，不给他们的雕像献花环和祭品。因为你们很清楚，同一动物被一些人视为神，被另一些人视为野兽，又被另一些人视为祭品。\par

% structure: p0051 source=h2 target=section reason=relative-heading-rank; base=h2

\section{第二十五章　基督徒弃绝假神}

其次，因为我们——曾从各族中崇拜塞墨勒之子巴克斯、拉托娜之子阿波罗（他们在与男人的爱情中做了甚至羞于提及之事），以及普洛塞庇娜和维纳斯（她们因对阿多尼斯的爱而疯狂，你们也庆祝她们的秘仪），或埃斯科拉庇俄斯，或那些被称为神祇中的某一位——如今藉耶稣基督学会了藐视这些，即使为此被处死也在所不惜，并将自己奉献给无生无受难的天主；我们确信祂从未被对安提俄珀或类似女人或伽倪墨得斯的情欲所驱使，也未被那位通过忒提斯得到帮助的百臂巨人解救，也未因此焦虑，为使其子阿喀琉斯因他的妾布里塞伊斯而毁灭许多希腊人。我们怜悯那些相信这些事的人，并知道那些发明这些事的是魔鬼。\par

% structure: p0053 source=h2 target=section reason=relative-heading-rank; base=h2

\section{第二十六章　基督徒不信任魔术师}

第三，因为在基督升天以后，魔鬼推出一些人，声称他们自己是神；你们不但不迫害他们，反而认为他们配得尊荣。有一个\Person{撒玛黎雅人}西满，来自一个名叫基托的村庄，在\Person{克劳狄乌斯·凯撒}统治时期，在你们罗马的皇城中，凭借魔鬼在他内运作的法术，行了强大的魔法。他被视为神，像神一样被你们以一座雕像尊崇，该雕像建立在\Place{台伯河}上，两桥之间，并带有罗马语言的铭文：—— \Quote{Simoni Deo Sancto,} \Quote{致神圣的西门神}。几乎所有\Person{撒玛黎雅人}，甚至其他民族中的少数人，都敬拜他，并承认他为第一位神；还有一个名叫\Person{赫肋纳}的女子，当时与他同行，原为妓女，他们说是他产生的第一个意念。还有一个名叫\Person{默南德}的人，也是\Person{撒玛黎雅人}，来自卡帕雷泰亚镇，是\Person{西满}的门徒，受魔鬼启示，据我们所知，他在\Place{安提约基雅}用他的法术欺骗了许多人。他说服那些跟随他的人，使他们永不会死，甚至现在还有一些活着的人持守他的观点。还有\Person{本都}人\Person{玛尔基昂}，至今仍在世，教导他的门徒相信某个比造物主更大的神。他藉魔鬼的帮助，使许多各族的人说亵渎的话，否认天主是这宇宙的创造者，并声称有另一位比他更大的存在做了更伟大的事。所有从这些人采纳意见者，如我们先前所说，被称为基督徒；正如那些不同意哲学家学说的人，却仍共同享有哲学家的名称。至于他们是否犯下那些神话般可耻的行为——翻倒灯、混乱交合、吃人肉——我们不知道；但我们确实知道，你们既不迫害他们，也不处死他们，至少不是因为他们的意见。但我已经有一部针对所有已存在的异端的论著，如果你们愿意阅读，我将把这部论著提供给你们。\par

% structure: p0055 source=h2 target=section reason=relative-heading-rank; base=h2

\section{第二十七章 遗弃婴孩之罪}

至于我们，我们受教导：遗弃初生婴孩是恶人的行为；我们受此教导，免得我们伤害任何人，也免得我们得罪天主，第一，因为我们看到几乎所有这样被遗弃的（不仅是女孩，也包括男孩）都被养大卖淫。正如古人据说饲养牛群、山羊群、绵羊群或牧马，现在我们看到你们养大孩子只为了这羞耻的用途；并且为了这污秽，在各民族中都可以找到大量的女性、\OriginalTerm{阴阳人}{hermaphrodites}以及行不可言喻的邪恶的人。你们从他们那里收取酬金，以及来自他们的税和课税，而你们本应从你们的领域中将他们消灭。任何人使用这样的人，除了不敬神、丑恶和不洁的交合，还可能是在与自己的孩子、亲戚或兄弟交合。有些人甚至卖淫自己的子女和妻子，有些人公开为\OriginalTerm{鸡奸}{sodomy}而阉割；他们将这些神秘仪式归于\Place{众神之母}，并与你们视为神的每一个一起画有一条蛇，一个伟大的象征和奥秘。事实上，你们公开并喝彩地做的事情，仿佛神圣之光被推翻和熄灭，你们却把这些归咎于我们；这实在对拒绝做这类事的我们毫无伤害，只伤害那些做这些事并作伪证控告我们的人。\par

% structure: p0057 source=h2 target=section reason=relative-heading-rank; base=h2

\section{第二十八章 天主对人的眷顾}

因为在我们中间，恶灵之首被称为蛇、\Person{撒殚}和魔鬼，你们可以通过查阅我们的著作得知。基督预言他必与他的部众以及跟随他的人一同被投入火中，受无尽期的惩罚。因为天主迟迟未行此事的原因，是祂对人类眷顾。祂预知有些人要通过悔改而得救，有些人甚至可能尚未出生。起初祂造人类具有思考和选择真理、行善的能力，使所有人在天主面前无可推诿；因为他们生来是理性而具有默想能力的。若有人不信天主眷顾这些事，他将因此暗示天主不存在，或者虽然存在却以恶为乐，或者像石头一样存在，而德行与恶行毫无意义，只是按人的意见视为善恶。这就是最大的亵渎和邪恶。\par

% structure: p0059 source=h2 target=section reason=relative-heading-rank; base=h2

\section{第二十九章 基督徒的贞洁}

再者[我们害怕遗弃孩子]，免得有些孩子不被拣起而死去，我们就成了杀人者。但无论我们结婚，仅为养育子女；或我们不结婚，我们保持贞洁。为使你们明白混乱交合不是我们的奥秘之一，我们中的一位不久前向\Place{亚历山大}总督\Person{斐理斯}递交了一份请愿书，恳请许可外科医生将他阉割。因为那里的外科医生说，没有总督的许可他们禁止这样做。当\Person{斐理斯}坚决拒绝签署这样的许可时，那青年保持单身，并满足于自己良心的认可，以及与他思想相同之人的认可。我们认为在此提及\Person{安提诺乌斯}并非不合时宜，他最近还活着，所有人都因恐惧而急切地敬拜他为神，尽管他们知道他是谁以及他的出身。\par

% structure: p0061 source=h2 target=section reason=relative-heading-rank; base=h2

\section{第三十章 基督难道不是术士吗？}

但恐怕有人会质问我们：有什么能阻止我们认为我们称为基督的那一位，一个由人所生的人，用法术行了他所谓的大能奇事，并由此显为天主子？我们现在要提供证据，不单凭断言，而是必然被那些预言[祂]的人说服，因为这些事发生之前已有预言，因为我们亲眼看见已发生和正在发生的事正如预言所说；我们认为这在你们看来也是最有力、最真实的证据。\par

% structure: p0063 source=h2 target=section reason=relative-heading-rank; base=h2

\section{第三十一章 论希伯来先知}

当时在犹太人中，有些人是天主的先知，藉着他们，预言之圣神预先宣告尚未发生而将要实现的事。他们的预言，当被说出和发出时，历代在犹太人中执政的诸王，按照先知们用希伯来语亲自编录成书的形式，都小心保存于自己手中。当埃及王\Person{托勒密}建立图书馆，并努力收集所有人的著作时，他也听说了这些先知，便派遣使者前往当时犹太人的王\Person{黑落德}那里，请求将先知的书卷送给他。\Person{黑落德}王果然将那些用上述希伯来语写成的书卷送去了。当埃及人发现其内容难以理解时，他又派人请求委任一些人将其翻译成希腊语。此事完成后，这些书卷便留在埃及人那里，直至如今。这些书卷也为全世界所有犹太人所拥有；但他们虽阅读，却不明白所说的，反而将我们视为仇敌；而且，如同你们一样，他们一有机会就杀害和惩罚我们，你们完全能够相信这一点。因为在最近爆发的犹太战争中，犹太人叛乱的领袖\Person{巴尔科赫巴}下令，唯独基督徒应被带至残酷刑罚，除非他们否认耶稣基督并口出亵渎。因此，在这些先知书中，我们发现了我们的基督耶稣被预言将要来临：由童贞女诞生，长大成人，治愈各种疾病和各样衰弱，复活死人，被憎恨，不被承认，被钉十字架，死亡，复活，升天，并且是也被称为天主子。我们也发现预言说，将有一些人由祂派遣到每个民族去宣扬这些事，并且人们要信靠祂，更是在外邦人中（而非在犹太人中）。在祂显现之前，祂已被预言，先是五千年前，再三千年前，然后两千年前，一千年前，又八百年前；因为在世代交替中，一位又一位先知相继兴起。\par

% structure: p0065 source=h2 target=section reason=relative-heading-rank; base=h2

\section{第三十二章 基督由\Person{梅瑟}预言}

当时，作为先知首位的\Person{梅瑟}，用这些话说道：\BibleQuote{权杖必不离开犹大，立法者不离他两足之间，直到那为祂所保留者来到；祂将是万民的期望，将祂的驴驹拴在葡萄树上，将祂的外袍在葡萄血中清洗。}\BibleRef{创 49:10} 你们可以仔细查考，确定直到何时犹太人有自己的立法者和君王。直到耶稣基督的时候，祂教导了我们，并解释了那些尚未被理解的预言，［他们仍有立法者］，正如神圣而属神的预言之圣神藉着\Person{梅瑟}所预言的：\BibleQuote{统治者不会离开犹太人，直到那为王权保留者来到}（因为犹大是犹太人的祖先，犹太人也是从他得名）；在祂（即基督）显现之后，你们开始统治犹太人，并获得了他们所有的领土。因为预言说：\BibleQuote{祂将是万民的期望}，意思是所有民族中都会有一些人期待祂再次来临。这一点你们确实可以亲眼看见，并因事实而信服。因为在所有人类种族中，都有一些人期待那位在犹太被钉十字架者，而祂被钉死后，那地立即作为战利品交给了你们。而预言说：\BibleQuote{将祂的驴驹拴在葡萄树上，将祂的外袍在葡萄血中清洗}，是基督将要遭遇和将要做之事的重要象征。因为有一匹驴驹拴在村庄入口处的葡萄树上，祂便吩咐祂的熟人当时把它牵来；牵来后，祂骑上去，进入耶路撒冷，那里有犹太人宏伟的圣殿，后来被你们摧毁。此后祂被钉十字架，使预言的其余部分得以应验。因为\BibleQuote{将祂的外袍在葡萄血中清洗}预言了祂将要受的苦难，藉祂的血洁净那些信靠祂的人。因为那被圣神藉着先知称为\BibleQuote{祂的外袍}的，就是那些信靠祂的人，在他们内住着天主的种子，即圣言。而称为\BibleQuote{葡萄血}的，表明那将要显现者会有血，但不是出于人的种子，而是出于天主的德能。在天主圣父和万有的主之后，首要的德能便是圣言，祂也是圣子；关于祂，我们将在下文讲述祂如何取了肉身并成为人。因为正如葡萄的血不是人而是天主所造，同样以此暗示，那血不是出于人的种子，而是出于天主的德能，正如我们上面所说。另一位先知\Person{依撒意亚}，用别的话语预言同样的事，这样说：\BibleQuote{由\Person{雅各伯}将升起一颗星，由\Person{叶瑟}的根将生出一朵花；万民将信赖祂的臂膀。}\BibleRef{依 11:1} 一颗光星已经升起，一朵花已从\Person{叶瑟}的根生出——这花就是基督。因为藉着天主的德能，祂由一位童贞女因\Person{雅各伯}的种子而受孕。\Person{雅各伯}是\Person{犹大}的父亲，我们已表明\Person{犹大}是犹太人的父亲；按照神谕，\Person{叶瑟}是祂的祖先，按血统祂是\Person{雅各伯}和\Person{犹大}的后裔。\par

% structure: p0067 source=h2 target=section reason=relative-heading-rank; base=h2

\section{第三十三章 基督诞生方式的预言}

再听依撒意亚如何明言预言祂将由童贞女所生；因为他这样说：\BibleQuote{看，童贞女将怀孕，生一个儿子，他们要称祂的名字为\InnerQuote{天主与我们同在}。} \BibleRef{依 7:14} 因为那些在人看来不可思议、似乎不可能的事，天主藉预言之神预言它们将要发生，以便当它们发生时，不致有不信，而是因着预言而有信德。但为避免有人不理解此刻所引的预言，反而用我们指责那些说朱庇特因情欲亲近妇女的诗人的同样理由来责备我们，让我们试解释这些话。因此，\BibleQuote{看，童贞女将怀孕}，意指童贞女将未经交合而怀孕。因为若她与任何人有过交合，她便不再是童贞女；但天主的德能临于童贞女，荫庇了她，使她虽是童贞女却怀孕了。当时奉命到那童贞女那里的天神来报佳音说：\BibleQuote{看，你将因圣神怀孕，生一个儿子，祂要称为至高者的儿子，你要给祂起名叫耶稣；因为祂要把自己的百姓从他们的罪恶中拯救出来，} \BibleRef{路 1:32} ——正如那些记录我们救主耶稣基督一切事迹的人所教导的，我们相信了他们，因为藉着我们现在引用的依撒意亚，预言之神也宣告了祂的诞生，正如我们之前所提的。因此，将圣神和天主的德能理解为圣言以外的任何事物都是错误的，圣言也是天主的首生子，正如前述先知梅瑟所宣告的；正是这圣言临于童贞女并荫庇她，使她怀孕，不是藉着交合，而是藉着德能。而耶稣这个名字在希伯来语中，在希腊语里意为\OriginalTerm{救主}{\GreekText{Σωτήρ}}。因此，天使也对童贞女说：\BibleQuote{你要给祂起名叫耶稣，因为祂要把自己的百姓从他们的罪恶中拯救出来。} 而且，先知们除了受天主圣言的感动之外，别无其他来源，这一点，我想连你们也会承认。\par

% structure: p0069 source=h2 target=section reason=relative-heading-rank; base=h2

\section{基督诞生地点的预言}

再听另一位先知米该亚预言祂将生于何地。他这样说：\BibleQuote{你，白冷，犹大地，在犹大的首领中并非最小；因为将由你出来一位领袖，祂将牧养我的百姓。} \BibleRef{米 5:2} 现在，在犹太地有一个村庄，离耶路撒冷三十五斯塔狄，耶稣基督便在那里诞生，你们也可以从在犹太地你们的第一任总督居里扭所作的人口登记册中查证。\par

% structure: p0071 source=h2 target=section reason=relative-heading-rank; base=h2

\section{其他应验的预言}

以及基督诞生后如何避开其他人直到长大成人，这事也已发生，请听关于此事的预言。有以下预言：——\BibleQuote{有一个婴孩为我们诞生了，有一个青年赐给了我们，政权必担在祂的肩上；} \BibleRef{依 9:6} 这象征十字架的权能，因为当祂被钉时，祂将祂的肩膀置于其上，这在后续的论述中将更加清楚。再者，同一位先知依撒意亚受预言之神感动，说：\BibleQuote{我整天向悖逆和抗拒的百姓伸出双手，他们走在不好的道路上。他们如今向我求审判，竟敢亲近天主。} \BibleRef{依 65:2} 又在其他话语中，通过另一位先知，祂说：\BibleQuote{他们刺穿了我的手和我的脚，为我的衣服拈阄。} 的确，说这些话的达味王和先知从未遭受这些事；但耶稣基督伸出双手，被犹太人钉在十字架上，他们反对祂，否认祂是基督。正如先知所说，他们折磨祂，把祂放在审判席上，说：审判我们吧。而\Quote{他们刺穿了我的手和我的脚}这句话，是指钉在祂手脚上的十字架钉子。祂被钉后，他们为祂的衣服拈阄，钉祂的人将衣服分了。这些事确实发生了，你们可以从\WorkTitle{本丢·彼拉多行传}中查证。我们还要引用另一位先知索福尼亚的预言，明确预言祂将骑驴驹进入耶路撒冷。话是这样说的：\BibleQuote{熙雍女子，你应欢喜踊跃；耶路撒冷女子，你应欢呼：看，你的君王来到你这里，谦逊地骑在一头驴上，一头驴驹上。} \BibleRef{匝 9:9}\par

% structure: p0073 source=h2 target=section reason=relative-heading-rank; base=h2

\section{预言的不同方式}

但当你们听到先知的话语仿佛是亲自说出时，你们不可认为是由被感动者自己说的，而是由感动他们的天主圣言说的。因为祂有时以预告未来的方式宣告将要发生的事；有时以天主、万有的主和父的位格说话；有时以基督的位格；有时以人民回应主或祂的父亲的位格，正如你们甚至可以在自己的作者中看到，一个人是整部作品的作者，却引入对话的人物。拥有先知书的犹太人并不明白这一点，因此基督来临时他们也不认识祂，反而恨我们这些说祂已来临并证明正如所预言祂被他们钉在十字架上的人。\par

% structure: p0075 source=h2 target=section reason=relative-heading-rank; base=h2

\section{圣父的发言}

为使你们也明白此事，曾从圣父的位格通过依撒意亚先知说了以下的话：\Quote{牛认识自己的主人，驴认识主人的槽；但以色列却不认识，我的子民也不明白。祸哉，犯罪的民族，充满罪恶的人民，邪恶的种子，悖逆的儿女，你们离弃了主。} 又在别处，当同一位先知以类似方式从圣父的位格发言时：\Quote{你们要为我建造什么样的殿宇？主说。天是我的宝座，地是我的脚凳。} \BibleRef{依 66:1} 又在另一处：\Quote{你们的月朔和安息日，我的灵魂憎恨；那大斋日和停工的日子，我不能容忍；你们前来朝见我，我也不听；你们的手沾满了血；你们献上细面、乳香，为我也是可憎的；羔羊的脂肪和公牛的血，我并不喜爱。因为谁向你们要求了这些？你们要松开邪恶的绳索，解开暴力的捆索，遮盖无家可归的人和赤身露体的人，把你的饼分给饥饿的人。} \BibleRef{依 1:14} 你们现在可以看出，通过先知从（天主的位格）教导的是什么样的事。\par

% structure: p0077 source=h2 target=section reason=relative-heading-rank; base=h2

\section{第三十八章 圣子的宣告}

当预言之神从基督的位格发言时，其宣告是这样：\Quote{我整天向悖逆顶嘴的百姓伸开双手，向那些行走不善之道的人。} \BibleRef{依 65:2} 又说：\Quote{我把我的背交给鞭打，我的颊交给掌击；我没有掩面躲避唾污的羞辱；主是我的助佑，因此我没有羞愧；我使我的脸像坚石，我知道我不会蒙羞，因为那使我成义的就在我身旁。} \BibleRef{依 50:6} 当祂说，\Quote{他们为我的衣服拈阄，穿透了我的手和脚。我躺下睡觉，又醒来，因为主支持了我。} 又说，\Quote{他们用嘴唇说话，摇头说：\InnerQuote{愿祂拯救自己。}} 你们可以知道，这一切事都发生在基督身上，由犹太人所为。因为当祂被钉十字架时，他们确实撇嘴摇头说：\Quote{让那使死人复活的拯救自己。} \BibleRef{玛 27:39}\par

% structure: p0079 source=h2 target=section reason=relative-heading-rank; base=h2

\section{第三十九章 圣神的直接预言}

当预言之神像预言将来要发生的事那样发言时，祂这样说：\Quote{因为法律将从熙雍而出，主的言语将从耶路撒冷而来。祂将在列国中施行审判，斥责许多民族；他们要把刀剑打成犁头，把枪矛打成镰刀；这民族不再举刀攻击那民族，也不再学习战争。} \BibleRef{依 2:3} 我们能使你们相信这事确实发生了。因为从耶路撒冷有十二个人出去到世上，他们目不识丁，无口才；但靠着天主的能力，他们向各族的人宣报，他们是基督派遣来教导众人天主圣言的；我们从前互相残杀，如今不仅不再向敌人作战，而且为不向审讯者说谎或欺骗，甘心为宣认基督而死。因为那句话\Quote{舌头已发誓，但心未发誓}，在此事上可以被我们效法。但如果你们所招募的士兵，宣了军誓，宁死也要忠于他们的誓言，胜过生命、父母、祖国和所有亲属，虽然你们不能给他们任何不朽之物，那么，我们这些热切渴望不朽的人，若不全心忍受一切，以获得那能赐予我们渴望之物的祂的恩赐，就实在是可笑的了。\par

% structure: p0081 source=h2 target=section reason=relative-heading-rank; base=h2

\section{第四十章 基督来临的预言}

又请听，关于那些传扬祂的教导并宣告祂显现的人是如何预言的。上述先知与君王藉着预言之神这样说：\Quote{日复一日发出言语，夜复一夜显示知识。没有言语，没有辞句，也听不见它们的声音。它们的声音传遍大地，它们的言语直达地极。祂在太阳中设置了祂的帐幕，祂如同新郎走出洞房，又如巨人般欢然奔跑。}此外，我们认为引用\Person{大卫}的其他一些预言也有必要并切合时宜；由此你们可以了解预言之神如何劝诫人们生活，以及祂如何预言了犹太人的君王\Person{黑落德}、犹太人本身，以及在他们中间的你们的总督\Person{比拉多}及其兵士所结成反对基督的阴谋；并预言祂如何被万民所信；天主怎样称祂为自己的儿子，并宣告要使祂的一切仇敌屈服于祂之下；以及魔鬼如何尽其所能地逃避天主父及万有的主宰的权力，也逃避基督本身的权力；并预言天主如何在审判之日来临之前召唤众人悔改。这些事这样陈述：\Quote{有福的人，不随从恶人的计谋，不站在罪人的道路上，不坐在傲慢者的座位上：他喜爱上主的法律；他昼夜默思祂的法律。他像一棵栽在溪水旁的树，按时结果，叶子不枯干，凡他所作的尽都顺利。恶人却不如此，他们像糠秕被风吹散。因此恶人在审判中站立不住，罪人在义人的会众中也如此。因为上主知道义人的道路；恶人的道路必要灭亡。列国为什么咆哮，万民为什么妄想？地上的君王起来，首领聚集商议，反对上主，反对祂的受傅者，说：让我们挣断他们的绳索，摆脱他们的轭。那坐在天上的必嗤笑他们，上主必讥讽他们。那时祂要在震怒中对他们说话，在烈怒中使他们惊慌。然而我已被祂立为君王，在\Place{熙雍}祂的圣山上，宣告上主的谕旨。上主对我说：你是我的儿子，我今日生了你。你向我求，我必将列国赐你为基业，将地极赐你为产业。你必用铁杖击打他们，将他们如同陶匠的器皿摔碎。所以，君王们，你们要明智；大地上的审判官们，要受教。以敬畏事奉上主，且要战兢喜乐。接受教诲，免得上主发怒，你们便从正道上灭亡，因为祂的怒气突然燃起。凡投靠祂的人都是有福的。}\par

% structure: p0083 source=h2 target=section reason=relative-heading-rank; base=h2

\section{第四十一章　预言被钉十字架}

又在另一预言中，预言之神通过同一位\Person{大卫}暗示，基督在被钉十字架之后将要为王，如下所述：\Quote{要向上主歌唱，全地，天天传扬祂的救恩。因为上主伟大，应受大赞美，祂当受敬畏超越众神。因为列邦的众神都是魔鬼的偶像，而上主创造了诸天。荣耀与赞美在祂面前，能力与荣耀在祂的圣所。要将荣耀归给上主，永远的父。领受恩宠，进入祂的临在，在祂的圣院内敬拜。愿全地在他面前战栗；愿它坚定，不动摇。愿列邦欢欣。上主从树上为王了。}\par

% structure: p0085 source=h2 target=section reason=relative-heading-rank; base=h2

\section{第四十二章　预言使用过去时态}

但当预言之神讲述将要发生的事，仿佛已经发生——正如在我们已经引用的段落中也可以观察到的那样——为了不使读者以此为借口（误解它们），我们也要使这一点非常清楚。祂绝对知道将要发生的事，就预言它们如同已经发生。你们若留意这些言语，就会明白这些话必须这样理解。以上引述的话，是\Person{大卫}在基督成为人并被钉十字架前一千五百年说的；在祂之前的人，以及祂的同时代人，没有一个是因被钉十字架而给外邦人带来喜乐的。但是我们的\Person{耶稣基督}，被钉十字架而死，又复活了，升天后为王；并且借着宗徒们以祂的名义在万民中所传扬的，给那些期待祂所应许的永生的人带来喜乐。\par

% structure: p0087 source=h2 target=section reason=relative-heading-rank; base=h2

\section{第四十三章　责任的主张}

但恐怕有人根据我们所说的话，以为我们主张一切发生的事皆由\OriginalTerm{宿命必然性}{fatal necessity}而发生，因为是预先知道而预言的，对此我们也作说明。我们从先知们学到并持守为真实的，就是惩罚、管束和善报都按照各人行为的功过而施予。因为若非如此，而是一切皆由命运决定，那么任何事都不在我们自己的能力之内。譬如，若命定此人为善、彼人为恶，则前者无功，后者无咎。再者，除非人类有凭\OriginalTerm{自由抉择}{free choice}避免恶而选择善的能力，否则无论他们行为如何，都不必负责任。但人确是由自由抉择而行善或跌倒，我们如此证明：我们看到同一个人能转到相反的事。若命定他必为善或为恶，他就绝不能兼具二者，也不能有如此多的转变。甚至也不会有些人善、有些人恶，因为那样我们就使命运成为恶的原因，显示命运自相矛盾；否则，就会出现先前已述的观点，即德行与恶行皆无意义，事物只因意见而被视为善恶。这，正如真道所显示的，是最大的不敬和邪恶。但我们断言，必然的命运乃是：选择善者获得相称的赏报，选择相反者获得应得的惩罚。因为天主造人不像其他事物，如树木和四足动物，不能凭选择行动；否则，人若不自己选择善，而是为此目的被造，就配不上赏报或赞美；同样，若他为恶，也配不上惩罚，因为他并非自己为恶，而是只能成为被造的样子。\par

% structure: p0089 source=h2 target=section reason=relative-heading-rank; base=h2

\section{第四十四章 预言不使其无效}

预言的圣神教导我们这事，借着梅瑟告诉我们，天主对最初受造的人这样说：\Quote{看，在你面前有善与恶：你要选择善。}又借着另一位先知依撒意亚，发出以下仿佛出自天主圣父和万有之主的言语：\Quote{你们要洗涤，自洁；从你们灵魂中除去邪恶；学习行善；为孤儿伸冤，为寡妇辩护；然后来，我们辩论，上主说：你们的罪虽似朱红，我将使它们白如羊毛；虽红如丹颜，我将使它们白如雪。你们若愿意并听从我，必享地之美物；但若不听从我，刀剑必吞噬你们：因为上主的口说了。}\BibleRef{依 1:16}而那话\Quote{刀剑必吞噬你们}，并非说不服从者必被刀剑杀死，而是天主的刀剑就是火，选择行恶的人成为其燃料。因此祂说：\Quote{刀剑必吞噬你们：因为上主的口说了。}若祂说的是切割并立即结束的刀剑，就不会说\Emph{吞噬}。同样，柏拉图说\Quote{选择者的罪责在他自己，天主无可责难}时，是从先知梅瑟那里取来而说的。因为梅瑟比所有希腊作家都更古老。无论哲学家还是诗人关于灵魂不死、死后惩罚、默观天上事物或诸如此类的学说，他们都是从先知们得到提示，因而能理解并解释这些事。因此，所有人中似乎都有真理的种子；但他们因主张矛盾之事而被指控未能准确理解（真理）。所以我们论及未来事件被预言时，并非说它们因宿命必然性而发生；但天主预知所有人将要行的一切事，并定意人的未来行为将按其各自价值得到报偿，祂便借着预言之神预言将照行为的功德给予相称的赏报，时常敦促人类努力并反省，显示祂眷顾并照顾人类。然而，因魔鬼的煽动，死刑被定下反对那些阅读\OriginalTerm{希斯塔斯普}{Hystaspes}、\OriginalTerm{西彼尔}{Sibyl}或先知们著作的人，为的是使他们因恐惧而阻止阅读者获得对善的认识，并使他们继续受制于自己；但这并非总能成功。因为不仅我们无畏地阅读这些书，而且如你们所见，还拿来供你们审阅，知道其内容将使众人喜悦。即使我们只说服少数人，我们的收获也很大；因为作为良善的农夫，我们将从主人那里得到赏报。\par

% structure: p0091 source=h2 target=section reason=relative-heading-rank; base=h2

\section{第四十五章 基督在天上坐于父右的预言}

万有的天主圣父使基督从死者中复活后，将带祂到天上，并保留祂在那里，直到祂征服祂的仇敌魔鬼，直到那些被祂预知为良善和德行完备者的数目满足，为此祂尚延缓终局——请听先知达味的话。他说：\Quote{上主对我主说：你坐在我右边，等我把你的仇敌放在你的脚下。上主必从耶路撒冷向你伸出权杖；你要在你仇敌中掌权。在你威能之日，政权与你同在，在你圣者的荣美中：我从早晨的怀中生了你。}祂所说的\Quote{祂必从耶路撒冷向你伸出权杖}，是预言那大能的圣言，即祂的宗徒们从耶路撒冷出发，四处宣讲的；尽管对于教导或哪怕承认基督之名的人，死刑已被判定，我们却到处接受并教导它。若你们也以敌意阅读这些话，你们所能做的，如前所述，不过是杀死我们；这确实对我们无害，却对你们和所有无理恨我们且不悔改的人，带来永火的惩罚。\par

% structure: p0093 source=h2 target=section reason=relative-heading-rank; base=h2

\section{第四十六章 基督降世之前的圣言}

但为免有人无理取闹，歪曲我们的教导，坚持说我们认为基督在齐林斯治下约一百五十年前诞生，后来在般雀·比拉多时代教导了我们所说的一切，并以此攻击我们，仿佛所有在祂之前出生的人都无需负责——让我们事先预见并解决这一难题。我们被教导说，基督是天主的首生者，并且我们已在上面阐明，祂是圣言，所有人类都分有祂；那些按理性生活的人，即使被视为无神论者，也是基督徒；例如在希腊人中，苏格拉底、赫拉克利特以及类似的人；在外邦人中，亚巴郎、阿纳尼雅、阿匝黎雅、米沙耳、厄里亚，以及许多其他人，他们的行为和名字我们此刻不再列举，因为知道这会冗长。所以，那些在基督之前生活且不按理性生活的人，是邪恶的，敌对基督，并杀害那些按理性生活的人。但是，那通过圣言之能、按照天主圣父和万有之主的旨意，由童贞女降生为人、被命名为耶稣、被钉十字架、死亡、复活并升天的人，一个明智的人将从已经充分讲述的内容中领悟。而我们，既然此刻对这一主题的证明不那么必要，将暂且转向那些紧迫之事的证明。\par

% structure: p0095 source=h2 target=section reason=relative-heading-rank; base=h2

\section{第四十七章 犹太地的荒芜已被预言}

那么，犹太之地将要荒芜，请听先知之神所言。这些话仿佛出自惊异于所发生之事的人民之口。它们是：\BibleQuote{熙雍成了旷野，耶路撒冷一片荒凉。我们圣所的家成了诅咒，我们祖先所祝福的荣耀被火烧毁，一切华美之物尽遭毁坏；祢对这些事缄默不言，又保持静默，使我们大大卑微。}\BibleRef{依 64:10}。你们确信耶路撒冷已如预言所说被荒废。关于它的荒凉，以及不许任何人居住其中，依撒意亚有以下预言：\BibleQuote{他们的土地荒凉，仇敌在他们面前吞吃它，他们中无人居住其中。}\BibleRef{依 1:7}。并且你们十分清楚，它被你们看守，不许任何人居住其中，且捕获入内的犹太人会被处死。\par

% structure: p0097 source=h2 target=section reason=relative-heading-rank; base=h2

\section{第四十八章 基督的作为与死亡已被预言}

关于我们的基督将要医治一切疾病并使死人复活，曾被预言，请听所说的话。有以下这些话：\BibleQuote{在祂来临时，瘸子要跳跃如鹿，哑巴的舌头要说话清晰；瞎子要看见，麻风病人要洁净；死人要复活，并行走。}\BibleRef{依 35:6}。并且祂确实行了这些事，你们可以从\WorkTitle{般雀·比拉多行传}中得知。至于先知之神如何预言祂和那些信靠祂的人将被杀害，请听依撒意亚所说的话。这些话是：\BibleQuote{看哪，义人灭亡，无人放在心上；公正之人被带走，无人思念。因邪恶面前，义人被带走，他的安息将是平安；他从我们中间被带走。}\BibleRef{依 57:1}。\par

% structure: p0099 source=h2 target=section reason=relative-heading-rank; base=h2

\section{第四十九章 祂被犹太人弃绝已被预言}

再者，同一位依撒意亚如何说，那不寻找祂的外邦民族将敬拜祂，而一直期盼祂的犹太人当祂来时却不认识祂。这些话是出自基督的口吻，它们是：\BibleQuote{我显现给那些不向我求问的人；我被那些不寻找我的人找到；我对那不呼求我名的民族说：我在这里。我向那悖逆顶嘴的百姓伸出双手，他们行不善之道，随从自己的罪恶；这百姓当着我的面惹我发怒。}\BibleRef{依 65:1}。因为犹太人拥有预言，一直期盼基督来临，却不认识祂；不仅如此，还羞辱地对待祂。但外邦人从未听过基督，直到宗徒们从耶路撒冷出发，向他们宣讲基督，并给予他们预言，他们才充满喜乐和信德，抛弃偶像，藉着基督将自己奉献给\OriginalTerm{不可生的天主}{Unbegotten God}。并且，针对那些承认基督的人将说出这些恶言，以及那些毁谤祂、说保持古老习俗是好事的人将陷入悲惨，这已被预知，请听依撒意亚简略所说的：\BibleQuote{祸哉，那些称甜为苦、称苦为甜的人。}\BibleRef{依 5:20}。\par

% structure: p0101 source=h2 target=section reason=relative-heading-rank; base=h2

\section{第五十章 祂的贬抑已被预言}

但是，祂为了我们而成为人，忍受了苦难和羞辱，并且祂将带着光荣再次来临，请听关于此事的预言；它们如下：\BibleQuote{因为他们把祂的灵魂交于死亡，祂被列于罪犯之中，祂承担了许多人的罪，并为悖逆者代求。看哪，我的仆人必行事明智，必被高举，必受极大的尊崇。许多人对你感到惊奇，你的容貌在人前如此毁损，你的荣耀如此隐藏；因此许多民族将惊奇，君王将因祂而闭口。因为那些未曾被告知关于祂的事，那些未曾听见的人将要明白。主啊，谁相信了我们的报告？主的手臂向谁启示了？我们在祂面前宣告如一个孩子，像旱地上的根。祂没有威仪，也没有光彩；我们看见祂时，也没有美丽可言：但祂的容貌比世人更被羞辱和毁损。一个受打击的人，知道如何忍受软弱，因为祂的面容被转开：祂被轻视，不被尊重。是祂承担了我们的罪，为我们受苦；我们却以为祂是受罚、被打击、被折磨。但祂为我们的过犯被刺透，为我们的罪恶被压伤；那带来和平的惩罚落在祂身上，因祂的创伤我们得了痊愈。我们都像羊一样迷了路，各人偏行己路。祂为我们的罪被交出；祂在一切苦难中不开口。祂被牵到宰杀之地，又像羊羔在剪毛的人前无声，祂不开口。在祂受辱时，祂的审判被夺去了。}\BibleRef{依 52:13} 因此，在祂被钉十字架后，连祂所有的相识都离弃了祂，否认了祂；后来，当祂从死者中复活并向他们显现，教导他们阅读那些预言，其中所有这一切都被预言将要发生，当他们看见祂升天，就相信了，并接受了从祂那里赐给他们的能力，就前往各民族，教导这些事情，并被称为宗徒。\par

% structure: p0103 source=h2 target=section reason=relative-heading-rank; base=h2

\section{第五十一章 基督的威严}

而且，预言之神为了向我们指明那遭受这些事者有着说不尽的根源，并统治祂的敌人，祂这样说：\BibleQuote{谁能述说祂的世代？因为祂的生命从地上被剪除：因他们的过犯祂来到死亡。我必将恶人归在祂的坟墓里，将富人归在祂的死中；因为祂没有行强暴，口中也没有诡诈。主喜悦将祂从鞭伤中洁净。若祂为罪而被交出，你的灵魂将看见祂的后裔延年益寿。主喜悦将祂的灵魂从痛苦中解救，显示祂光明，用知识塑造祂，使那为多人丰盛服务的义人成义。祂将承担我们的罪过。因此，祂将承受许多，并分得强者的掠物；因为祂的灵魂被交付于死亡；祂被列于罪犯之中；祂承担了许多人的罪，并因他们的过犯而被交出。}\BibleRef{依 53:8} 请听，根据预言，祂是如何升天的。这是这样说的：\BibleQuote{举起天门的门闩，打开吧，好让光荣的君王进来。谁是这位光荣的君王？主，大能而有力的。} 还有，祂将如何带着光荣再次从天上降临，请听耶肋米亚先知关于此事所说的话。祂的话是：\BibleQuote{看，祂如同人子乘着天上的云彩而来，祂的天使们与祂一同前来。}\BibleRef{达 7:13}\par

% structure: p0105 source=h2 target=section reason=relative-heading-rank; base=h2

\section{第五十二章 预言的确定应验}

既然我们证明已经发生的一切事情在它们发生之前都已由先知预言过，那么我们也必须相信那些同样被预言、但尚未发生的事，必定会发生。因为正如已经发生的事在预言时即已应验，即使不被知晓，同样，那些剩余的事，即使不被知晓、不被相信，也必将应验。因为先知们宣告了祂的两次来临：一次是已经过去的，当祂作为一位受羞辱和受苦的人来到；第二次则是按预言，祂将带着光荣从天上降临，由祂的天使军队陪同，那时祂将复活所有曾经活过的人的身体，将配得者的身体穿上不朽，而将恶人的身体，赋予永恒的知觉，与邪恶的魔鬼一同投入永火之中。这些事也被预言将要发生，我们将证明。厄则克耳先知曾说：\BibleQuote{关节与关节相连，骨与骨相接，肌肉重新生长；一切膝都要跪拜上主，一切口舌都要承认祂。}\BibleRef{则 37:7} 至于恶人将受怎样的知觉和刑罚，请听与此类似的预言如下：\BibleQuote{他们的虫子不死，他们的火不灭；}\BibleRef{依 66:24} 到那时他们将悔改，但已无济于事。至于犹太人民当他们看见祂在光荣中降临时将说的话和做的事，匝加利亚先知如此预言：\BibleQuote{我将命令四风收集四散的子女；我将命令北风带来他们，南风不可阻拦。然后在耶路撒冷将有极大的哀悼，不是口头或嘴唇的哀悼，而是心灵的哀悼；他们将撕裂的不是衣服，而是心灵。各支派各自哀悼，然后他们将仰望他们所刺透的那一位；他们要说：主啊，为什么祢使我们偏离祢的道路？我们祖先所祝福的光荣，为我们变成了羞辱。}\par

% structure: p0107 source=h2 target=section reason=relative-heading-rank; base=h2

\section{第五十三章 预言摘要}

虽然我们可以提出许多其他预言，但我们止步于此，认为这些足以说服那些有耳可听、有耳可悟的人；同时，我们也考虑到那些人能够看出，我们并非仅仅断言而没有证据，如同那些关于所谓朱庇特之子的传说一样。因为，我们有什么理由相信一个被钉十字架的人是非受生天主的首生子，并且祂将要审判全人类，除非我们在祂降生为人之前就找到了关于祂的预言，并且我们看到事情已然实现——犹太人的土地荒凉，以及各族的人通过宗徒被祂的教导说服，弃绝了他们旧日的习俗，在他们的迷惑中行事；是的，我们看到我们自己，也知道来自外邦人的基督徒比来自犹太人和撒玛黎雅人的更多、更真实？因为所有其他人类种族都被预言之神称为外邦人；但犹太和撒玛黎雅种族被称为以色列的支派和雅各伯的家。至于预言中预表说从外邦人中相信的人要比从犹太人和撒玛黎雅人中更多，我们将引述如下：它这样说：\BibleQuote{不生育、不生产的女子啊，欢乐吧！未经产痛的女子，欢呼高唱吧！因为被弃者的子女比有夫者的子女更多。}\BibleRef{依 54:1} 因为所有外邦人都\Quote{荒凉}脱离了真天主，事奉他们手所做的；而犹太人和撒玛黎雅人虽由先知领受了天主的话，并一直期待基督，但祂来的时候，除了少数人以外，他们并未认出祂；关于这少数人，预言之神曾通过依撒意亚预言他们将被拯救。祂以他们的口吻说：\BibleQuote{若非主给我们留下种籽，我们早已如同索多玛和哈摩辣了。}\BibleRef{依 1:9} 因为据梅瑟记载，索多玛和哈摩辣是不虔敬之人的城市，天主用火与硫磺焚烧，使之倾覆，城中居民无一得救，只有一个外来的加色丁人，名叫罗特，与他的女儿们被救出。关心此事的人至今仍可看到那整个地区荒凉、被焚烧，并且寸草不生。为了显示从外邦人中如何被预言为更真实、更相信，我们将引述依撒意亚先知的话；他说如下：\BibleQuote{以色列人心未受割礼，外邦人却肉体未受割礼。}因此，诸如此类的事物，若亲眼看见，足以使那些接受真理、不固执己见、不被私欲支配的人产生信服和信仰。\par

% structure: p0109 source=h2 target=section reason=relative-heading-rank; base=h2

\section{第五十四章 异教神话的起源}

但是，那些传授诗人编造的神话的人，对学习它们的年轻人提不出任何证据；我们要证明这些神话是由邪恶的魔鬼所影响而编造的，目的是欺骗和误导人类。因为魔鬼们通过先知听到基督将要来临，以及人类中的不虔敬者将被火惩罚，他们就推出许多人称为朱庇特之子，以为这样能使人们以为关于基督所说的事仅仅是奇谈怪论，如同诗人所说的那些事一样。这些事既在希腊人中，也在所有民族中流传（魔鬼们在那里听到先知预言基督将特别被相信）；但他们在听先知所说的话时并未准确理解，而是模仿了我们基督所说的，如同犯错的人一样。我们将阐明这一点。先知梅瑟，正如我们已说过的，比所有作家都古老；通过他，也如先前所说，曾如此预言：\BibleQuote{权杖不离犹大，律法不离他脚间，直到那应得权杖者来到，万民都归顺他；祂将驴驹拴在葡萄树上，将祂的衣裳在葡萄血中漂洗。}\BibleRef{创 49:10} 因此，魔鬼们听到这些预言的话，就说巴克斯是朱庇特的儿子，并宣称他是葡萄的发现者，他们将酒（或驴）列入他的秘仪；他们教导说，他被人撕碎后升了天。因为梅瑟的预言中没有明确暗示要来的那位是天主子，也不知道祂骑着驴驹是留在地上还是升天，而\Quote{驴驹}这个词既可以指驴的驹子，也可以指马的驹子，他们不知道预言的那位会带来驴驹还是马驹作为祂来临的记号，也不知道祂是天主子（正如我们所说）还是人子，他们就宣称佩伽索斯马的骑手柏勒洛丰——一个凡人所生的人——自己骑着他的马升天了。当他们从另一位先知依撒意亚听说祂将由童贞女所生，并亲自升天时，他们就宣称珀耳修斯就是所说的那位。当他们知道所引用的预言所说\BibleQuote{强壮如巨人奔跑他的路程}时，他们就宣称赫拉克勒斯强壮，并走遍了全地。当他们又得知曾预言祂将医治各种疾病、使死人复活时，他们就制造了埃斯科拉庇俄斯。\par

% structure: p0111 source=h2 target=section reason=relative-heading-rank; base=h2

\section{第五十五章 十字架的象征}

但在任何情况下，即使是在那些被称为朱庇特之子的当中，他们也没有模仿被钉十字架的事；因为他们不明白它，关于它的一切都是象征性地陈述的。而这一件事，正如先知所预言的，是他权能与角色的最大象征；这由我们观察到的事物也得到证明。因为请考虑世上万物，没有这种形状，它们能否被管理或具有任何共性？因为海洋无法航行，除非那个被称为帆的胜利标记安然留在船上；没有它，土地无法耕种：挖掘者和工匠无法做他们的工作，除非使用具有这种形状的工具。而人的形体与无理性的动物不同，仅仅在于它是直立的，双手伸展，脸上有从额头延伸出来的所谓鼻子，生物通过它呼吸；而这显示的没有别的形状，只有十字架的形状。所以先知曾说：\BibleQuote{我们面前的呼吸是主基督。}而这一形状的权能由你们自己的象征物——所谓\Quote{\OriginalTerm{军旗}{vexilla}}和胜利纪念物——所显示，你们所有的国家财产都使用这些，用作你们权力和政府的徽章，尽管你们不自觉。并且用这种形状，你们在皇帝死后祝圣他们的像，并通过铭文称他们为神。因此，既然我们既通过理性也通过明显的形状劝勉了你们，并尽了我们最大的能力，我们知道现在即使你们不信，我们也是无可指责的；因为我们的本分已经完成并结束了。\par

% structure: p0113 source=h2 target=section reason=relative-heading-rank; base=h2

\section{第五十六章 魔鬼仍迷惑人}

但是那些邪灵并不满足于在基督显现之前说那些据说是朱庇特之子的人是由他而生；但当祂显现并生于人间，当他们得知祂如何被先知所预言，并知道祂应被万民信仰和期待之后，他们又像前面所说的那样，推出了另外一些人，撒玛黎雅人西满和默南德尔，他们藉魔法行了许多大能的事，欺骗了许多人，并且至今仍使他们受骗。因为即使在你们中间，正如我们之前所说，西满在克劳狄皇帝统治时期在皇城罗马，使神圣元老院和罗马人民大为惊异，以至于他被视为神，并像你们尊为神的其他人一样，被以雕像尊崇。因此我们祈求，神圣元老院和你们的人民，可以同你们一起，成为我们这个陈情的仲裁者，以便如果有人被那个人的学说所纠缠，他可以认识真理，从而能够逃脱错误；至于那座雕像，如果你们愿意，就毁掉它。\par

% structure: p0115 source=h2 target=section reason=relative-heading-rank; base=h2

\section{第五十七章 并引起迫害}

魔鬼也不能说服人认为不会有惩罚恶人的烈火；正如他们不能使基督在祂来了之后被隐藏一样。但他们只能做到这一点：那些过无理性生活、在邪恶习俗中放荡长大、并固执己见的人，会杀害和仇恨我们；我们不但不仇恨他们，而且如事实证明，我们怜悯他们，努力引导他们悔改。因为我们不怕死亡，因为大家都承认我们必定要死；没有新事，但在这种事物管理下，万事依然相同；如果那些即使享受这些事物一年的人感到厌倦，他们应当留意我们的教义，以便他们能永远地、无痛苦无匮乏地生活。但如果他们相信死后一无所有，并宣称死者进入无知觉状态，那么当他们使我们脱离此生的痛苦和必需时，他们就成了我们的恩人，并且证明自己是邪恶、不人道和固执的。因为他们杀我们并非为了释放我们，而是断绝我们，使我们被剥夺生命和快乐。\par

% structure: p0117 source=h2 target=section reason=relative-heading-rank; base=h2

\section{第五十八章 并兴起异端者}

而且，如我们之前所说，魔鬼推出了本都的马西翁，他甚至至今仍在教导人们否认天主是天上地上万物的创造者，否认先知所预言的基督是天主之子，并宣扬在万有的创造者之外还有另一位神，以及另一位儿子。许多人相信了这个人，好像只有他知道真理，并且嘲笑我们，尽管他们对自己所说的没有证据，却被无理地如同狼叼羊一样被带走，成为无神论教义和魔鬼的掠物。因为那些被称为魔鬼的，别无企图，只是要引诱人离开造他们的天主和祂的首生者基督；那些不能超升于地之上的人，他们已将他们钉牢，并且现在仍将他们钉牢于尘世的事物和手中的作品之上；但那些致力于默观神性事物的人，他们暗中击退；如果他们缺乏智慧的清醒、纯洁而无激情的生活，他们就将他们驱入无神状态。\par

% structure: p0119 source=h2 target=section reason=relative-heading-rank; base=h2

\section{第五十九章 柏拉图受惠于梅瑟}

为使你们知晓，柏拉图是从我们的教师——我们是指通过先知传递的记载——那里借用了他的说法，即天主改变了无形状的质料，创造了世界，请听通过梅瑟所说的话，正如上面所示，梅瑟是最早的先知，比希腊作家更为古老；通过他，预言的神灵指示了天主起初如何以及由什么材料形成世界，如此说：\BibleQuote{在起初，天主创造了天地。大地是无形空虚的，深渊上是一片黑暗，天主的神在水面上运行。天主说：\InnerQuote{有光！}就有了光。}因此，不仅柏拉图以及赞同他的人，而且我们自己也学到了，你们也可以被说服：整个世界是藉天主的话由梅瑟先前所说的物质所造成的。而诗人们称为\OriginalTerm{厄瑞玻斯}{Erebus}的东西，我们知道梅瑟早就说过。\BibleRef{申 32:22}\par

% structure: p0121 source=h2 target=section reason=relative-heading-rank; base=h2

\section{第六十章 柏拉图关于十字架的教导}

关于天主子在柏拉图《\WorkTitle{蒂迈欧篇}》中的生理学讨论，他说：\Quote{他将他十字形地放在宇宙中}，他同样是从梅瑟那里借用的；因为在梅瑟的著作中记载了那时，当以色列子民走出埃及，在旷野中时，他们遇到了毒蛇、蝮蛇和各类毒蛇，咬死了百姓；梅瑟受天主的感动和影响，取了铜，制成了十字架的形像，放在圣所中，并对百姓说：\BibleQuote{你们若注视这形像，并相信，就必借此得救。}见：\BibleRef{户 21:8} 这事完成后，记载说蛇就死了，并且传说百姓就这样逃脱了死亡。柏拉图读了这些事，没有准确理解，也没有领悟这是十字架的形像，反而以为是十字形放置，就说仅次于第一位天主的\OriginalTerm{能力}{power}被十字形地放在宇宙中。至于他提到第三位，是因为他读了我们上面所说的梅瑟的话：\BibleQuote{天主的圣神运行在水面上。}他给予第二位是那与天主同在的\OriginalTerm{圣言}{Logos}，他说他被十字形地放在宇宙中；第三位是那被说成运行在水面上的圣神，他说：\Quote{第三位围绕第三者。}请听预言之神如何通过梅瑟预示将有大火。他这样说：\BibleQuote{永火将降下，吞灭到深渊。}见：\BibleRef{申 32:22} 所以我们并非持有与他人相同的意见，而是所有人都在模仿我们的教导。在我们这里，这些事情可以从那些甚至不认识字母形状、言语粗俗野蛮、但心思明智而有信仰的人那里听到和学习；有些人甚至残疾瞎眼；这样你们可以明白，这些事不是出于人的智慧，而是出于天主的能力所说出的。\par

% structure: p0123 source=h2 target=section reason=relative-heading-rank; base=h2

\section{第六十一章 基督徒的圣洗}

我也要说明我们藉基督成为新人后，如何将自己奉献给天主；免得我们省略此事，在我们所作的说明中显得不公。凡被说服并相信我们所教导和所说的是真实的，并承诺能够照此生活的人，都被教导要祈祷、以斋戒恳求天主，求赦免他们过去的罪，我们也与他们一同祈祷和斋戒。然后，他们被我们带到有水的地方，并按照我们自己所重生的方式重生。因为，因着天主、宇宙之父和主的名，以及我们的救主耶稣基督的名，和圣神的名，他们接受水的洗涤。因为基督也说：\BibleQuote{除非你们重生，否则不能进入天国。}见：\BibleRef{若 3:5} 一旦出生，再进入母腹是不可能的，这对所有人都是明摆着的。至于那些犯罪而悔改的人如何逃脱他们的罪，如我上面所写，依撒意亚先知已宣告；他这样说：\BibleQuote{你们要洗濯、自洁，从我眼前除去你们的恶行；要学习行善，寻求正义，责打压迫者，为孤儿伸冤，为寡妇辩护。然后来，我们辩论，上主说。你们的罪虽像朱红，必变成雪白；虽红如丹颜，必白如羊毛。你们若乐意服从，将享用地上的美物；若不听从，反而悖逆，刀剑必吞灭你们：因为上主的口说了。}见：\BibleRef{依 1:16}\par

关于这一\OriginalTerm{仪式}{rite}，我们从宗徒那里得知了理由。由于我们出生时，并非出于自己的知识或选择，而是由于父母的结合，并在恶习和邪恶的教养中长大；为了使我们不再成为必然和无知的孩子，而成为选择和知识的孩子，并在水中获得对先前所犯罪过的赦免，于是，对于那选择重生并已悔改自己罪过的人，要呼号天主、宇宙之父和主的名；引导那将要受洗的人到水盆的人只以这名呼叫他。因为没有人能说出那不可言说的天主的名；若有人胆敢说有名字，他是在妄念中疯狂。这洗涤被称为\OriginalTerm{光照}{illumination}，因为学习这些事的人在理智上得到光照。并在那在般雀·比辣多手下被钉十字架的耶稣基督的名，以及那通过先知预言了关于耶稣的一切的圣神的名下，受光照的人得以洗涤。\par

% structure: p0126 source=h2 target=section reason=relative-heading-rank; base=h2

\section{第六十二章 恶魔对它的模仿}

事实上，魔鬼们听到先知宣告了这一洗涤，便煽动进入他们神庙的人，在接近他们时用奠酒和燔祭，也给自己洒水；他们还让他们在离开祭献后、进入设有偶像的圣所之前，彻底洗净自己。至于司祭们给进入神庙敬拜的人的命令，要他们脱鞋，魔鬼们得知了上面提到的先知梅瑟所发生的事，也模仿了这些事。因为在当时，当梅瑟奉命下到埃及，领出在那里的以色列子民时，他在阿拉伯地牧放他舅父的羊群，我们的基督以火从荆棘丛中显现与他说话，并说：\BibleQuote{脱下你的鞋，走近来听。}他脱下鞋走近后，听到他必须下到埃及，领出在那里的以色列子民；他从以火显现与他说话的基督那里领受了大力，便下去领出了百姓，行了伟大奇妙的事；如果你们想知道，你们可以从他的著作中准确了解。\par

% structure: p0128 source=h2 target=section reason=relative-heading-rank; base=h2

\section{第六十三章 天主如何显现给梅瑟}

所有犹太人至今仍教导说，无名的天主曾对梅瑟说话；因此预言之神通过上面提到的先知依撒意亚指责他们说：\BibleQuote{牛认识自己的主人，驴也认识主人的槽，但以色列却不认识我，我的百姓也不明白。}\BibleRef{依 1:3}同样，耶稣基督也指责犹太人，因为他们不知道父是什么，子是什么；祂亲自说：\BibleQuote{除了父外，没有人认识子；除了子外，也没有人认识父，以及子所启示的人。}\BibleRef{玛 11:27}圣言就是天主子，正如我们之前所说。祂被称为天使和宗徒，因为祂宣告我们应当知道的一切，并被派遣宣告所启示的一切；正如我们主亲自说：\BibleQuote{听从我的，就是听从那派遣我的。}\BibleRef{路 10:16}从梅瑟的著作中也可以看出这一点；因为其中写道：\BibleQuote{天主的天使在荆棘丛的火焰中对梅瑟说：\InnerQuote{我是自有者，我是亚巴郎的天主，依撒格的天主，雅各伯的天主，你父亲的天主；下到埃及去，把我的百姓领出来。}}\BibleRef{出 3:6}如果你想了解后续，可以从同一著作中得知；这里无法全部叙述。但所写的这些足以证明耶稣基督是天主子和祂的宗徒，祂从太初就是圣言，有时以火的形状显现，有时以天使的形像；但如今，因天主的旨意，为人类成为了人，忍受了魔鬼煽动无知的犹太人加给祂的一切苦难；犹太人虽然清楚地看到梅瑟著作中写着：\BibleQuote{天主的天使在荆棘的火焰中对梅瑟说：\InnerQuote{我是自有者，我是亚巴郎的天主，依撒格的天主，雅各伯的天主。}}却坚持说这话的是宇宙之父和创造者。因此预言之神斥责他们说：\BibleQuote{以色列不认识我，我的百姓不明白我。}\BibleRef{依 1:3}再者，耶稣如我们已指出的，当祂与他们在一起时，说：\BibleQuote{除了父外，没有人认识子；除了子外，也没有人认识父，以及子所愿意启示的人。}\BibleRef{玛 11:27}因此，犹太人一直认为对梅瑟说话的是宇宙之父，虽然那对梅瑟说话的确实是天主子，祂被称为天使和宗徒，他们被预言之神和基督本人公正地指责为既不知道父，也不知道子。因为那些声称子就是父的人，被证明既不认识父，也不知道宇宙之父有一个子；这子也是天主的首生之言，就是天主。古时祂以火的形状和天使的形像向梅瑟和其他先知显现；但如今在你们统治的时期，祂如我们之前所说，按照父的旨意，由童贞女成为人，为拯救那些相信祂的人，祂忍受了藐视和苦难，为通过死亡和复活战胜死亡。那在荆棘中对梅瑟所说的话：\BibleQuote{我是自有者，我是亚巴郎的天主，依撒格的天主，雅各伯的天主，你父亲的天主。}\BibleRef{出 3:6}这意味着他们虽然死了，却仍然存在，并且是基督自己的人。因为他们是所有人中最早寻求天主的；亚巴郎是依撒格之父，依撒格是雅各伯之父，正如梅瑟所记载。\par

% structure: p0130 source=h2 target=section reason=relative-heading-rank; base=h2

\section{第六十四章 对真理的进一步曲解}

从前面所说，你可以明白魔鬼如何模仿梅瑟的话，声称\OriginalTerm{普洛塞庇娜}{Proserpine}是\OriginalTerm{朱庇特}{Jupiter}的女儿，并煽动人们在泉源旁以\OriginalTerm{科瑞（少女或女儿）}{Kore (Cora)}之名设立她的雕像。因为，正如我们上面所写，梅瑟说：\BibleQuote{起初天主创造了天地。大地混沌空虚，深渊上一片黑暗，天主的神在水面上运行。}因此，他们模仿这里所说的天主的神在水面上运行，就声称普洛塞庇娜（或科瑞）是朱庇特的女儿。同样，他们还狡猾地虚构\OriginalTerm{密涅瓦}{Minerva}是朱庇特的女儿，并非通过两性结合，而是知道天主藉圣言构思并创造了世界，他们说密涅瓦是最初的概念\OriginalTerm{概念}{\GreekText{ἔννοια}}，我们认为这非常荒谬，竟以女性形象来表现概念。同样，那些被称为朱庇特之子的其他人的行为足以谴责他们。\par

% structure: p0132 source=h2 target=section reason=relative-heading-rank; base=h2

\section{第六十五章 圣事的施行}

但我们在这样洗净了那被说服并赞同我们教导的人之后，带他到那些被称为弟兄的人聚集的地方，以便我们共同为自己、为\OriginalTerm{受洗者（受光照者）}{baptized [illuminated] person}以及为各处所有其他人献上恳切的祈祷，使我们学习真理后，也得以因行为而被视为好公民和诫命的遵守者，从而获得永久的救恩。祈祷结束后，我们以亲吻彼此问候。然后，将饼和一杯掺水的酒带到弟兄们的主持者面前；他接过这些东西，藉圣子和圣神的名赞美光荣宇宙之父，并长时间感恩，因为我们得以从祂手中领受这些恩赐。当他结束祈祷和感恩时，所有在场的人以说\Quote{阿们}表示同意。这个阿们一词在希伯来语中相当于\OriginalTerm{但愿如此}{\GreekText{γένοιτο}}。当主持者感恩完毕，所有的人都表示同意后，那些被我们称为执事的人，将经过感恩的饼和掺水的酒分给每一位在场的人领受，并给缺席的人带一份。\par

% structure: p0134 source=h2 target=section reason=relative-heading-rank; base=h2

\section{第六十六章 论圣体圣事}

此食物在我们中称为\OriginalTerm{圣体}{\GreekText{Εὐχαριστία}}，只有相信我们所教真理、且已受那为使罪得赦并重生的圣洗、并依照基督所命生活的人，才可领受。因为我们领受这些，并非如同普通的饼和普通的酒，而是正如我们的救主耶稣基督藉着天主的圣言降生成人，为我们的救恩具有血肉，同样，我们也被告知：那因祂的言语祈祷而受祝祷、并藉转化滋养我们血与肉的食物，就是那位成为肉身的耶稣的血与肉。因为宗徒们在他们所撰写的回忆录（即所谓\WorkTitle{福音}）中，如此传授了赐给他们的命令：耶稣拿起饼来，祝谢后说：\Quote{你们应这样做来纪念我，\BibleRef{路 22:19}这是我的身体；}同样，他拿起杯来祝谢后说：\Quote{这是我的血；}并只赐给了他们。邪恶的魔鬼在密特拉的神秘仪式中模仿了这事，命令做同样的事。因为，在入教者的秘礼中，面包和一杯水与某些咒语一同放置，此事你们或是知道，或是可以了解。\par

% structure: p0136 source=h2 target=section reason=relative-heading-rank; base=h2

\section{第六十七章 基督徒的每周敬礼}

此后我们不断互相提醒这些事。我们中的富裕者帮助贫穷者；我们常相聚；对于所供给的一切，我们藉着祂的圣子耶稣基督并藉着圣神赞颂万有的创造者。在称为主日的那一天，凡住在城里或乡下的人都聚集到一处，宗徒们的回忆录或先知们的著作被宣读，时间许可多久便读多久；然后，当读者停止后，主持者口头训导，并劝勉效法这些善行。随后我们一同起立祈祷，正如我们先前所说的，祈祷结束后，便奉上饼、酒和水，主持者照自己的能力献上祈祷和感恩，民众应声说\Quote{阿们}；于是分给每人，并共享那已祝谢的恩物，也为缺席者由执事送一份。那些富裕而自愿的人，各按己意捐献；收集的款项交由主持者保管，他照顾孤儿、寡妇、因病或其他原因贫困的人、被囚的人，以及我们中寄居的旅客，总之，照顾一切有需要的人。但主日是我们举行公共集会的日子，因为那是天主在黑暗与物质中施行变化而创造世界的第一日；并且我们的救主耶稣基督在同一日从死者中复活。因为他在星期六的前一天被钉十字架；并在星期六的后一天，即太阳日，显现给祂的宗徒和门徒，教导了他们这些事，我们也将其提交给你们，供你们思量。\par

% structure: p0138 source=h2 target=section reason=relative-heading-rank; base=h2

\section{第六十八章 结语}

如果这些事在你们看来合理而真实，请予尊重；但若显得荒谬，就当作胡说而鄙视，不要对那些没有行恶的人如同对待敌人一样判处死刑。因为我们预先警告你们，若你们继续不义，必不能逃脱天主将来的审判；而我们自己将邀请你们做天主所喜悦的事。尽管凭最伟大、最杰出的皇帝\Person{哈德良}、你们的父亲的书信，我们可以要求你们按我们所愿下令审判，但我们作出这一申诉和解释，并非基于哈德良的决定，而是因为知道我们所求是公义的。我们附上了哈德良书信的副本，使你们知道我们在此事上说真话。以下是其副本：——\par

% structure: p0140 source=h2 target=section reason=relative-heading-rank; base=h2

\section{附录}

% structure: p0141 source=h3 target=subsection reason=relative-heading-rank; base=h2

\subsection{哈德良皇帝为基督徒辩护的书信}

我已收到你的前任、最杰出的人物\Person{塞雷尼乌斯·格拉尼亚努斯}写给我的信；我不愿对此信默不作声，以免无辜者受扰，并给告密者提供作恶的机会。因此，如果贵省的居民愿意支持这一请愿，在某个法庭上控告基督徒，我不禁止他们这样做。但我不会允许他们仅仅使用恳求和喊叫。因为若有人想要提出控告，你对此作出判决，这更为公正。因此，若有人提出控告，并提供证据证明这些人做了任何违反法律的事，你应依罪行大小判处刑罚。而且，凭赫丘利，你应特别注意：若有人仅以诬蔑控告这些人中的任何一个，你应按其邪恶程度更严厉地惩罚他。\par

% structure: p0143 source=h3 target=subsection reason=relative-heading-rank; base=h2

\subsection{安东尼努斯皇帝致亚细亚共同议会的书信}

皇帝凯撒\Person{提图斯·埃利乌斯·哈德良·安东尼·奥古斯都·庇护}，最高大司祭，在其护民官任期第十五年，第三次担任执政官，祖国之父，致亚洲行省公会，谨此致意：我本以为众神自会留意，不让此类罪犯逃脱。因为如果他们有能力，他们自己定会更愿惩罚那些拒绝敬拜他们的人；然而正是你们给这些人带来麻烦，将他们所持的信仰控告为无神论者的意见，并指控他们其他一些我们无法证实的事情。但对他们而言，被人以为是为其所受的指控而死倒是有利的；他们宁可慷慨舍弃自己的生命，也不愿顺从你们所要求的服从，以此战胜你们。至于那些已经发生和正在发生的地震，你们向我们提起它们是不恰当的；每当它们发生，你们就灰心丧气，从而将自己的行为与这些人的行为形成对比；因为他们在天主面前的信心远比你们大得多。的确，你们在这样的时刻似乎无视众神，忽略圣殿，毫不承认对天主的敬拜。因此你们嫉妒那些侍奉他的人们，并将他们迫害至死。关于此类人，其他一些行省的总督也曾写信给我最神圣的父亲；他回复说，他们不应打扰此类人，除非发现他们企图反对罗马政府。至于我本人，许多人已就此类人向我报告，我也遵循我父亲的判断予以答复。但如果有人仅仅以此类人的身份对其中任何一人提出控告，就让被告被宣告无罪，即使他被发现是这样一个；但让原告承担法律责任。\par

% structure: p0145 source=h3 target=subsection reason=relative-heading-rank; base=h2

\subsection{\Person{马尔库斯·奥勒留}致元老院的书信，其中他证明基督徒是他胜利的原因}

皇帝凯撒\Person{马尔库斯·奥勒留·安东尼}，日耳曼尼库斯、帕提库斯、萨尔马提库斯，致罗马人民与神圣元老院，谨此致意：我向你们解释了我的伟大计划，以及在日耳曼边境我历经艰辛所获得的利益，这是因为我被敌人包围；我本人被困在\Place{卡农图姆}，由七十四支步兵队包围，离敌军九英里。敌军就在眼前，侦察兵向我们指出，我们的将军\Person{庞培安努斯}向我们显示，我们附近有一支九十七万七千人的混杂人群，我们确实看见了；我被这支大军包围，身边只有第一、第十、双倍和海军军团组成的一个营。在审视了自己的处境和我的军队，对比野蛮人和敌人的庞大数量后，我迅速转向向祖国的众神祈祷。但他们无视了我，于是我召唤了我们中间称为基督徒的人。经过询问，我发现他们人数众多、队伍庞大，便对他们发怒，这实属不当；因为后来我了解了他们的力量。于是他们开始战斗，不是准备武器、军械或号角；因为这样的准备为他们所憎恶，出于他们心中所怀有的天主。因此，很可能那些我们认为是无神论者的人，有天主作为他们的统治力量扎根于他们的良心之中。因为他们伏在地上，不仅为我祈祷，也为所有在场的军队祈祷，好让他们摆脱眼前的干渴和饥荒。因为五天来我们得不到水，因为没有水；我们正处在日耳曼腹地，在敌人的领土上。而就在他们伏在地上向天主（一位我不认识的天主）祈祷的同时，天上降下水来，对我们来说是极为清凉的甘露，但对罗马的敌人却是一场毁灭性的冰雹。我们立刻认识到，随着祈祷，天主的临在——一位不可战胜、不可毁灭的天主——降临了。因此，基于这一点，让我们宽恕那些是基督徒的人，免得他们为我们祈祷并获得这样的武器来对付我们自己。我建议，不得因某人是基督徒而控告他。但如果有人被发现指控基督徒为基督徒，我希望证明：被控为基督徒并承认自己是基督徒的人，除这一点外没有被指控任何其他事；但起诉他的人将被活活烧死。我还希望，负责行省治理的人不得强迫承认并证实此事的基督徒收回其信仰，也不得囚禁他。我希望这些事由元老院的一项法令予以确认。我命令将此敕令公布于\Place{图拉真广场}，以便人阅读。总督\Person{维特拉西乌斯·波利奥}将负责将其传达至所有邻近行省，并确保任何希望使用或持有它的人都能从我现公布的文件中无障碍地获得一份副本。\par

\SourceRecord{Translated by Marcus Dods and George Reith.；Ante-Nicene Fathers；Vol. 1.；Edited by Alexander Roberts, James Donaldson, and A. Cleveland Coxe.；Buffalo, NY: Christian Literature Publishing Co.,；1885.；<http://www.newadvent.org/fathers/0126.htm>.}
